% Generated by the blueprint tool (blueprint-pdf/1) from a blueprint snapshot; edit the
% Markdown sources, not this file.  Compile with `latexmk -pdf` or by running
% `pdflatex` twice (the second run settles the contents and the page references).
% Images, if any, are read from assets/ next to this file.
\documentclass[11pt]{article}
\usepackage[T1]{fontenc}
\usepackage[utf8]{inputenc}
\usepackage{lmodern}
\usepackage{textcomp}
\usepackage[a4paper,margin=2.5cm,headheight=14pt]{geometry}
\usepackage{amsmath,amssymb,amsthm}
\usepackage{amscd}  % KaTeX's CD diagrams; part of amsmath, so always there
\usepackage{mathtools}
\usepackage{bm}
\usepackage{graphicx}
\usepackage{xcolor}
\usepackage{array,booktabs,tabularx}
\usepackage{fancyvrb}
\usepackage[normalem]{ulem}
\usepackage{microtype}
% Optional, and only with their fonts: a small TeX installation has the
% packages without the fonts, which fails at the first formula.
\IfFileExists{mathrsfs.sty}{\IfFileExists{rsfs10.tfm}{\usepackage{mathrsfs}}{}}{}
\IfFileExists{stmaryrd.sty}{\IfFileExists{stmary10.tfm}{\usepackage{stmaryrd}}{}}{}
\IfFileExists{cancel.sty}{\usepackage{cancel}}{}
\usepackage[hyphens]{url}
\usepackage[pdfusetitle,bookmarksnumbered=false]{hyperref}
\usepackage{bookmark}
\definecolor{bplink}{HTML}{1F4E9C}
\definecolor{bpmuted}{HTML}{5F676E}
\hypersetup{colorlinks=true,linkcolor=bplink,urlcolor=bplink,citecolor=bplink}
\setlength{\parindent}{0pt}
\setlength{\parskip}{0.55em plus 0.2em minus 0.1em}
% Lists as tight as the page's.
\makeatletter
\def\@listi{\leftmargin\leftmargini\parsep 0pt\topsep 0.3em plus 0.1em\itemsep 0.1em plus 0.1em}
\let\@listI\@listi
\def\@listii{\leftmargin\leftmarginii\labelwidth\leftmarginii\advance\labelwidth-\labelsep\parsep 0pt\topsep 0.1em\itemsep 0.1em}
% An ordered list that does not start at 1.
\newcommand{\bpstartat}[1]{\setcounter{\@enumctr}{\numexpr#1-1\relax}}
\makeatother
\setcounter{tocdepth}{3}
\emergencystretch=3em
\allowdisplaybreaks

% A numbered statement, set as a paper sets one: "Theorem 2.1 (Title)."
% #1 is the lead ("Theorem~2.1"), #2 what follows it in roman.
\newenvironment{bpstatement}[2]{\par\addvspace{0.6\baselineskip}\noindent\textbf{#1}#2\textbf{.}\enspace\ignorespaces}
  {\par\addvspace{0.4\baselineskip}}
\newenvironment{bpstatementit}[2]{\begin{bpstatement}{#1}{#2}\itshape\ignorespaces}{\end{bpstatement}}
% A step: a relation with prose of its own, read under the object it starts from.
\newenvironment{bpstep}[1]{\par\addvspace{0.3\baselineskip}\begin{list}{}{\leftmargin=1.6em\rightmargin=0pt\topsep=0pt}\item\relax
  {\small\color{bpmuted}#1}\par\nopagebreak}{\end{list}\par}
% Notes about where an object is written out.
\newenvironment{bpnote}{\par\small\color{bpmuted}\itshape}{\par}
\newcommand{\bpomitted}[1]{\par{\small\color{bpmuted}\itshape #1\par}}
\newcommand{\bpbodyheading}[1]{\par\addvspace{0.4\baselineskip}\noindent\textbf{#1}\par\nopagebreak}
\newcommand{\bprule}{\par\medskip\noindent\hfil\rule{0.3\linewidth}{0.4pt}\hfil\par\medskip}
\newcommand{\bpboundary}[1]{\par\nopagebreak{\small\color{bpmuted}#1\par}\nopagebreak}
% A reference to something this export leaves out, and one that resolves nowhere.
\newcommand{\bpout}[1]{#1\textsuperscript{\hyperlink{bp-outside}{\ensuremath{\ast}}}}
\newcommand{\bpbroken}[1]{\textcolor{red!70!black}{[#1?]}}
\newcommand{\bpslug}[1]{\texttt{#1}}
% A bibliography entry: label, anchor, entry; the label hangs in the margin.
\newcommand{\bpbibitem}[3]{\par\setbox0\hbox{[#1]\enspace}\ifdim\wd0<5em\setbox0\hbox to 5em{[#1]\hfil}\fi
  \hangindent=\wd0 \hangafter=1 \noindent\box0\phantomsection\label{#2}#3\par\smallskip}
% A character the export knows no way to set.
\newcommand{\bpglyph}[1]{\fbox{\scriptsize #1}}
% Images, at their own size unless they are wider than the line.
\newsavebox{\bpimagebox}
\newcommand{\bpimage}[1]{\sbox{\bpimagebox}{\includegraphics{#1}}%
  \ifdim\wd\bpimagebox>\linewidth\resizebox{\linewidth}{!}{\usebox{\bpimagebox}}\else\usebox{\bpimagebox}\fi}
\newcommand{\bpmissingimage}[1]{\fbox{\parbox{0.8\linewidth}{\centering\small\color{bpmuted}Image not included: #1}}}
\DefineVerbatimEnvironment{bpverbatim}{Verbatim}{fontsize=\small,xleftmargin=1em}
\renewcommand{\qedsymbol}{\ensuremath{\blacksquare}}
\hypersetup{pdftitle={Weil Conjectures — 16 The Fourier transform},pdfsubject={Chapter 16 The Fourier transform; all levels},pdfcreator={blueprint}}
% KaTeX commands the website knows
\providecommand{\R}{\mathbb{R}}\providecommand{\N}{\mathbb{N}}\providecommand{\Z}{\mathbb{Z}}
\providecommand{\Q}{\mathbb{Q}}\providecommand{\C}{\mathbb{C}}
\providecommand{\Reals}{\mathbb{R}}\providecommand{\reals}{\mathbb{R}}\providecommand{\Complex}{\mathbb{C}}
\providecommand{\cnums}{\mathbb{C}}\providecommand{\natnums}{\mathbb{N}}\providecommand{\Bbb}[1]{\mathbb{#1}}
\providecommand{\bold}[1]{\mathbf{#1}}\providecommand{\frak}[1]{\mathfrak{#1}}\providecommand{\mathbbm}[1]{\mathbb{#1}}
\providecommand{\lang}{\langle}\providecommand{\rang}{\rangle}\providecommand{\bra}[1]{\langle #1\rvert}
\providecommand{\ket}[1]{\lvert #1\rangle}\providecommand{\braket}[1]{\langle #1\rangle}
\providecommand{\set}[1]{\{#1\}}\providecommand{\Set}[1]{\left\{#1\right\}}
\providecommand{\infin}{\infty}\providecommand{\isin}{\in}\providecommand{\plusmn}{\pm}\providecommand{\sub}{\subset}
\providecommand{\sube}{\subseteq}\providecommand{\supe}{\supseteq}\providecommand{\exist}{\exists}
\providecommand{\bull}{\bullet}\providecommand{\larr}{\leftarrow}\providecommand{\rarr}{\rightarrow}
\providecommand{\lrarr}{\leftrightarrow}\providecommand{\Larr}{\Leftarrow}\providecommand{\Rarr}{\Rightarrow}
\providecommand{\Lrarr}{\Leftrightarrow}\providecommand{\uarr}{\uparrow}\providecommand{\darr}{\downarrow}
\providecommand{\harr}{\leftrightarrow}\providecommand{\lArr}{\Leftarrow}\providecommand{\rArr}{\Rightarrow}
\providecommand{\hArr}{\Leftrightarrow}\providecommand{\alef}{\aleph}\providecommand{\alefsym}{\aleph}
\providecommand{\real}{\Re}\providecommand{\image}{\Im}\providecommand{\weierp}{\wp}
\providecommand{\Alpha}{\mathrm{A}}\providecommand{\Beta}{\mathrm{B}}\providecommand{\Epsilon}{\mathrm{E}}
\providecommand{\Zeta}{\mathrm{Z}}\providecommand{\Eta}{\mathrm{H}}\providecommand{\Iota}{\mathrm{I}}
\providecommand{\Kappa}{\mathrm{K}}\providecommand{\Mu}{\mathrm{M}}\providecommand{\Nu}{\mathrm{N}}
\providecommand{\Omicron}{\mathrm{O}}\providecommand{\Rho}{\mathrm{P}}\providecommand{\Tau}{\mathrm{T}}
\providecommand{\Chi}{\mathrm{X}}\providecommand{\omicron}{o}\providecommand{\gt}{>}\providecommand{\lt}{<}
\providecommand{\And}{\&}\providecommand{\hdashline}{\hline}\providecommand{\operatornamewithlimits}[1]{\operatorname*{#1}}
\providecommand{\llbracket}{[\![}\providecommand{\rrbracket}{]\!]}\providecommand{\cancel}[1]{#1}
\providecommand{\mathscr}{\mathcal}
% the characters beyond ASCII this document uses
\DeclareUnicodeCharacter{2161}{\textup{II}}
\DeclareUnicodeCharacter{2162}{\textup{III}}
% the project's macros ([katex.macros] in blueprint.toml)
\providecommand{\A}{}\renewcommand{\A}{\mathbf A}
\providecommand{\AEn}{}\renewcommand{\AEn}{A_{E,n}}
\providecommand{\Aaf}{}\renewcommand{\Aaf}{\mathbf A}
\providecommand{\Ab}{}\renewcommand{\Ab}{\mathsf{Ab}}
\providecommand{\AffSch}{}\renewcommand{\AffSch}{\operatorname{AffSch}}
\providecommand{\Aone}{}\renewcommand{\Aone}{\mathbf A^1}
\providecommand{\Bl}{}\renewcommand{\Bl}{\operatorname{Bl}}
\providecommand{\Br}{}\renewcommand{\Br}{\mathrm{Br}}
\providecommand{\C}{}\renewcommand{\C}{\mathbb C}
\providecommand{\Cech}{}\renewcommand{\Cech}{\check{H}}
\providecommand{\Coext}{}\renewcommand{\Coext}{\operatorname{Coext}}
\providecommand{\Cone}{}\renewcommand{\Cone}{\operatorname{Cone}}
\providecommand{\Cons}{}\renewcommand{\Cons}{\operatorname{Cons}}
\providecommand{\Cref}{}\renewcommand{\Cref}[1]{\text{#1}}
\providecommand{\D}{}\renewcommand{\D}{D}
\providecommand{\DLAFT}{}\renewcommand{\DLAFT}{\operatorname{FT}}
\providecommand{\DLAFr}{}\renewcommand{\DLAFr}{\operatorname{Fr}}
\providecommand{\DLAQ}{}\renewcommand{\DLAQ}{{\overline{\mathbf Q}_\ell}}
\providecommand{\DLRQ}{}\renewcommand{\DLRQ}{{\overline{\mathbf Q}_\ell}}
\providecommand{\Dbc}{}\renewcommand{\Dbc}{D^b_c}
\providecommand{\Dc}{}\renewcommand{\Dc}{D_{\mathrm{cons}}}
\providecommand{\Dctf}{}\renewcommand{\Dctf}{D^b_{ctf}}
\providecommand{\Dic}{}\renewcommand{\Dic}{D_{\IndCons}}
\providecommand{\Dplus}{}\renewcommand{\Dplus}{D^{+}}
\providecommand{\End}{}\renewcommand{\End}{\operatorname{End}}
\providecommand{\Et}{}\renewcommand{\Et}{\mathrm{\acute{e}t}}
\providecommand{\Ex}{}\renewcommand{\Ex}{\mathrm{Ex}}
\providecommand{\Ext}{}\renewcommand{\Ext}{\operatorname{Ext}}
\providecommand{\F}{}\renewcommand{\F}{\mathbb F}
\providecommand{\FT}{}\renewcommand{\FT}{\operatorname{FT}}
\providecommand{\Fbar}{}\renewcommand{\Fbar}{{\overline{\mathbf F}_q}}
\providecommand{\Fell}{}\renewcommand{\Fell}{\mathbf F_\ell}
\providecommand{\Fib}{}\renewcommand{\Fib}{\operatorname{Fib}}
\providecommand{\Fix}{}\renewcommand{\Fix}{\operatorname{Fix}}
\providecommand{\Fp}{}\renewcommand{\Fp}{{\mathbf F_{p}}}
\providecommand{\Fq}{}\renewcommand{\Fq}{{\mathbf F_q}}
\providecommand{\Fqn}{}\renewcommand{\Fqn}{{\mathbf F_{q^{n}}}}
\providecommand{\Frob}{}\renewcommand{\Frob}{\operatorname{Frob}}
\providecommand{\GL}{}\renewcommand{\GL}{\operatorname{GL}}
\providecommand{\Gal}{}\renewcommand{\Gal}{\operatorname{Gal}}
\providecommand{\Hens}{}\renewcommand{\Hens}{\operatorname{Hens}}
\providecommand{\Hilb}{}\renewcommand{\Hilb}{\operatorname{Hilb}}
\providecommand{\Hom}{}\renewcommand{\Hom}{\operatorname{Hom}}
\providecommand{\IndCons}{}\renewcommand{\IndCons}{\operatorname{IndCons}}
\providecommand{\Kar}{}\renewcommand{\Kar}{\operatorname{Kar}}
\providecommand{\LT}{}\renewcommand{\LT}{\operatorname{LT}}
\providecommand{\Lie}{}\renewcommand{\Lie}{\operatorname{Lie}}
\providecommand{\Loc}{}\renewcommand{\Loc}{\operatorname{Loc}}
\providecommand{\Mod}{}\renewcommand{\Mod}{\mathsf{Mod}}
\providecommand{\N}{}\renewcommand{\N}{\mathbb{N}}
\providecommand{\Needspace}{}\renewcommand{\Needspace}{\par}
\providecommand{\O}{}\renewcommand{\O}{{\mathcal{O}}}
\providecommand{\Ob}{}\renewcommand{\Ob}{\mathrm{Ob}}
\providecommand{\Oe}{}\renewcommand{\Oe}{\mathcal O_E}
\providecommand{\PP}{}\renewcommand{\PP}{\mathbf{P}}
\providecommand{\Perf}{}\renewcommand{\Perf}{\operatorname{Perf}}
\providecommand{\Pic}{}\renewcommand{\Pic}{\mathrm{Pic}}
\providecommand{\Pm}{}\renewcommand{\Pm}{\mathbf P^{m}}
\providecommand{\Pone}{}\renewcommand{\Pone}{\mathbf P^{1}}
\providecommand{\Pqn}{}\renewcommand{\Pqn}{\mathbf P^{n}}
\providecommand{\ProFin}{}\renewcommand{\ProFin}{\operatorname{ProFin}}
\providecommand{\Proet}{}\renewcommand{\Proet}{\mathrm{pro\acute{e}t}}
\providecommand{\Q}{}\renewcommand{\Q}{\mathbb Q}
\providecommand{\Qb}{}\renewcommand{\Qb}{\overline{\mathbf Q}_{\ell}}
\providecommand{\Qbar}{}\renewcommand{\Qbar}{\overline{\mathbf Q}}
\providecommand{\QbarX}{}\renewcommand{\QbarX}{\overline{\mathbf Q}_{\ell,X}}
\providecommand{\Qbarl}{}\renewcommand{\Qbarl}{{\overline{\mathbf Q}_{\ell}}}
\providecommand{\Qell}{}\renewcommand{\Qell}{{\mathbf{Q}_{\ell}}}
\providecommand{\Ql}{}\renewcommand{\Ql}{\mathbf Q_{\ell}}
\providecommand{\Qlb}{}\renewcommand{\Qlb}{{\overline{\mathbf Q}_\ell}}
\providecommand{\Qlbar}{}\renewcommand{\Qlbar}{\overline{\mathbf Q}_\ell}
\providecommand{\R}{}\renewcommand{\R}{\mathbb{R}}
\providecommand{\RG}{}\renewcommand{\RG}{\operatorname{\mathrm{R}\Gamma}}
\providecommand{\RGamma}{}\renewcommand{\RGamma}{R\Gamma}
\providecommand{\RGammac}{}\renewcommand{\RGammac}{R\Gamma_c}
\providecommand{\RHom}{}\renewcommand{\RHom}{R\mathcal{H}om}
\providecommand{\RHomX}{}\renewcommand{\RHomX}{R\mathcal H om_X}
\providecommand{\Rat}{}\renewcommand{\Rat}{\operatorname{IsRat}}
\providecommand{\Res}{}\renewcommand{\Res}{\operatorname{Res}}
\providecommand{\Rlim}{}\renewcommand{\Rlim}{R\!\varprojlim}
\providecommand{\SIm}{}\renewcommand{\SIm}{\operatorname{SIm}}
\providecommand{\Sch}{}\renewcommand{\Sch}{\mathsf{Sch}}
\providecommand{\Sh}{}\renewcommand{\Sh}{\operatorname{Sh}}
\providecommand{\Span}{}\renewcommand{\Span}{\operatorname{span}}
\providecommand{\Spec}{}\renewcommand{\Spec}{\operatorname{Spec}}
\providecommand{\Tr}{}\renewcommand{\Tr}{\operatorname{Tr}}
\providecommand{\Trsm}{}\renewcommand{\Trsm}{\operatorname{Tr}^{\mathrm{sm}}}
\providecommand{\Wt}{}\renewcommand{\Wt}{\operatorname{Wt}}
\providecommand{\Z}{}\renewcommand{\Z}{\mathbb{Z}}
\providecommand{\Zar}{}\renewcommand{\Zar}{\mathrm{zar}}
\providecommand{\Zell}{}\renewcommand{\Zell}{{\mathbf{Z}_{\ell}}}
\providecommand{\Zl}{}\renewcommand{\Zl}{{\mathbf Z_\ell}}
\providecommand{\affet}{}\renewcommand{\affet}[1]{{#1}^{\mathrm{aff}}_{\mathrm{ét}}}
\providecommand{\affproet}{}\renewcommand{\affproet}[1]{{#1}^{\mathrm{aff}}_{\mathrm{proét}}}
\providecommand{\boxt}{}\renewcommand{\boxt}{\boxtimes_A^{\mathrm L}}
\providecommand{\cF}{}\renewcommand{\cF}{\mathcal{F}}
\providecommand{\cHom}{}\renewcommand{\cHom}{R\mathcal{H}\!om}
\providecommand{\cI}{}\renewcommand{\cI}{\mathcal{I}}
\providecommand{\cL}{}\renewcommand{\cL}{\mathcal L}
\providecommand{\cM}{}\renewcommand{\cM}{\mathcal M}
\providecommand{\cN}{}\renewcommand{\cN}{\mathcal N}
\providecommand{\cO}{}\renewcommand{\cO}{\mathcal O}
\providecommand{\cT}{}\renewcommand{\cT}{\mathcal T}
\providecommand{\checked}{}\renewcommand{\checked}{\color{red}\aftergroup\restorecolor}
\providecommand{\coeff}{}\renewcommand{\coeff}{\operatorname{coeff}}
\providecommand{\coker}{}\renewcommand{\coker}{\operatorname{coker}}
\providecommand{\colim}{}\renewcommand{\colim}{\operatorname{colim}}
\providecommand{\constructedby}{}\renewcommand{\constructedby}[1]{\par\smallskip\noindent\textcolor{purple!60!black}{\textsf{Constructed by:} }
  \footnotesize\sloppy\nolinkurl{#1}\par}
\providecommand{\deepresult}{}\renewcommand{\deepresult}[1]{\par\smallskip\noindent\textcolor{orange!45!black}{
  \textsf{Blueprint role: named deep theorem node; see the cluster proof architecture.}}
  \par\noindent\textcolor{orange!45!black}{\textsf{Mathematical source:} }#1\par}
\providecommand{\dt}{}\renewcommand{\dt}{\mathbin{\otimes^{\mathbf L}}}
\providecommand{\dw}{}\renewcommand{\dw}{\operatorname{dw}}
\providecommand{\et}{}\renewcommand{\et}[1]{{#1}_{\mathrm{ét}}}
\providecommand{\ev}{}\renewcommand{\ev}{\operatorname{ev}}
\providecommand{\fib}{}\renewcommand{\fib}{\operatorname{fib}}
\providecommand{\formalizationnote}{}\renewcommand{\formalizationnote}[1]{\begin{quote}\small\raggedright\textbf{Formalization note.} #1\end{quote}}
\providecommand{\fpqc}{}\renewcommand{\fpqc}{X_{\mathrm{fpqc}}}
\providecommand{\gammacont}{}\renewcommand{\gammacont}{\Gamma_{\mathrm{cont}}}
\providecommand{\gm}{}\renewcommand{\gm}{\mathbf G_m}
\providecommand{\hocolim}{}\renewcommand{\hocolim}{\operatorname{hocolim}}
\providecommand{\holim}{}\renewcommand{\holim}{\operatorname{holim}}
\providecommand{\hrB}{}\renewcommand{\hrB}{\@ifnextchar]{\hrEnd}{\hrStep}}
\providecommand{\hrEnd}{}\renewcommand{\hrEnd}{\@firstoftwo{}}
\providecommand{\hrStep}{}\renewcommand{\hrStep}{\@firstoftwo{\hrB}}
\providecommand{\id}{}\renewcommand{\id}{\operatorname{id}}
\providecommand{\kunFl}{}\renewcommand{\kunFl}{\mathbf F_\ell}
\providecommand{\kunZn}{}\renewcommand{\kunZn}{\mathbf Z/n\mathbf Z}
\providecommand{\kuncd}{}\renewcommand{\kuncd}{\operatorname{cd}}
\providecommand{\kundim}{}\renewcommand{\kundim}{\dim}
\providecommand{\kunotimes}{}\renewcommand{\kunotimes}[1]{\otimes^{\mathbf L}_{#1}}
\providecommand{\kunsupp}{}\renewcommand{\kunsupp}{\operatorname{Supp}}
\providecommand{\lotimes}{}\renewcommand{\lotimes}{\otimes^{\mathrm L}}
\providecommand{\m}{}\renewcommand{\m}{{\mathfrak{m}}}
\providecommand{\n}{}\renewcommand{\n}{{\mathfrak{n}}}
\providecommand{\nagataGap}{}\renewcommand{\nagataGap}[1]{\par\noindent\textbf{Formalization gap.} #1\par}
\providecommand{\openobligation}{}\renewcommand{\openobligation}[1]{\par\smallskip\noindent\textcolor{red!65!black}{\textsf{Open obligation:} }
  \footnotesize\sloppy\nolinkurl{#1}\par}
\providecommand{\p}{}\renewcommand{\p}{{\mathfrak{p}}}
\providecommand{\pr}{}\renewcommand{\pr}{\operatorname{pr}}
\providecommand{\preshv}{}\renewcommand{\preshv}[1]{\mathrm{PreShv}(#1)}
\providecommand{\pro}{}\renewcommand{\pro}[1]{\mathrm{Pro}(#1)}
\providecommand{\proet}{}\renewcommand{\proet}{X_{\mathrm{proét}}}
\providecommand{\q}{}\renewcommand{\q}{{\mathfrak{q}}}
\providecommand{\rad}{}\renewcommand{\rad}{\operatorname{rad}}
\providecommand{\rank}{}\renewcommand{\rank}{\operatorname{rank}}
\providecommand{\red}{}\renewcommand{\red}{\color{red}}
\providecommand{\restorecolor}{}\renewcommand{\restorecolor}{\color{black}}
\providecommand{\retiredobligation}{}\renewcommand{\retiredobligation}[1]{\par\smallskip\noindent\textcolor{red!75!black}{\textsf{Retired invalid route:} }
  \footnotesize\sloppy\nolinkurl{#1}\par}
\providecommand{\rk}{}\renewcommand{\rk}{\operatorname{rk}}
\providecommand{\sep}{}\renewcommand{\sep}{\textnormal{sep}}
\providecommand{\shv}{}\renewcommand{\shv}[1]{\mathrm{Shv}(#1)}
\providecommand{\spec}{}\renewcommand{\spec}[1]{\mathrm{Spec}(#1)}
\providecommand{\stacks}{}\renewcommand{\stacks}[1]{\href{https://stacks.math.columbia.edu/tag/#1}{\textnormal{Stacks, Tag~#1}}}
\providecommand{\stacksproject}{}\renewcommand{\stacksproject}[1]{{\cite[\href{https://stacks.math.columbia.edu/tag/#1}{Tag #1}]{stacks-project}}}
\providecommand{\statusnote}{}\renewcommand{\statusnote}[1]{\begin{quote}\small\raggedright\textbf{Status.} #1\end{quote}}
\providecommand{\tildeideal}{}\renewcommand{\tildeideal}[2]{\widetilde{(#1, #2)}}
\providecommand{\tildering}{}\renewcommand{\tildering}[3]{{#1}_{\widetilde{(#2, #3)}}}
\providecommand{\tr}{}\renewcommand{\tr}{\operatorname{tr}}
\providecommand{\wt}{}\renewcommand{\wt}{\operatorname{wt}}

\begin{document}
\pagestyle{plain}
\begin{center}
{\LARGE\bfseries Weil Conjectures\par}
\vspace{0.7em}
{\large Chapter 16: The Fourier transform\par}
\vspace{0.5em}
{\small\color{bpmuted}Snapshot 2026-09-26, revision 8e4cb88e2a\par}
\end{center}
\vspace{0.5em}
{\small\tableofcontents}
\bigskip\hrule\bigskip

% object: chap-The-Fourier-transform (blueprint/chap-The-Fourier-transform/_section.md)
\section*{\texorpdfstring{16\quad The Fourier transform}{16 The Fourier transform}}
\phantomsection\label{bp1}%
\addcontentsline{toc}{section}{\texorpdfstring{16\quad The Fourier transform}{16 The Fourier transform}}
% text: prose-16-01
\phantomsection\label{bp2}%
Theorem \bpout{14.26} controls the Frobenius eigenvalues of a real sheaf at the points of a curve. The bound we need next, however, is of a different kind. Let \(M\) be an irreducible lisse sheaf, pure of weight \(w\), on a dense open \(j:U_0\hookrightarrow\mathbf A^1\). The curve case of the Deligne–Laumon theorem comes down to showing that \(H^1_c(\mathbf A^1,j_*M)\) has weights at most \(w+1\). This is a cohomology group, not the stalk of a sheaf at a point, and Rankin's argument does not reach it. Is there a sheaf on another curve whose stalk at one point is this group, and whose weight we can compute from that of \(M\)? Up to a shift, the Fourier transform of \((j_*M)[1]\) is such a sheaf.

For functions on \(\mathbf F_q\) the discrete Fourier transform already does what we want. Fix a nontrivial additive character \(\psi\) and put \(\hat t(y)=\sum_xt(x)\psi(xy)\). Two properties matter. First, the value \(\hat t(0)\) is the sum of all values of \(t\), so a sum over the line becomes a single value. Second, Plancherel's formula \(\sum_y|\hat t(y)|^2=q\sum_x|t(x)|^2\) holds, by orthogonality of characters.

To see Plancherel compute an absolute value, let \(\chi\) be a nontrivial character of \(\mathbf F_q^\times\), extended by \(\chi(0)=0\), and let \(g=\sum_x\chi(x)\psi(x)\) be its Gauss sum. Substituting \(x\mapsto x/y\) shows that \(\hat\chi(y)=\overline\chi(y)g\) for \(y\ne0\), so the transform has the same absolute value \(|g|\) at all \(q-1\) nonzero points, while \(\hat\chi(0)=0\). Plancherel therefore reads \((q-1)|g|^2=q(q-1)\), and we find \(|g|=\sqrt q\). In the language of weights, this means that the transform carries a function of weight \(0\) to one of weight \(1\). For the quadratic character of \(\mathbf F_3\), moreover, Example \hyperref[bp8]{16.2} finds this Gauss sum, up to sign, as a Frobenius eigenvalue on a cohomology group.

The proof in Chapter \bpout{17} follows the same two steps for \(M\): first we show that the transform is pure, and then Plancherel finds its weight.

Grothendieck's function–sheaf dictionary (Proposition \bpout{12.28}) turns the transform into a functor. Let \(\mathcal K_\psi=\mu^*\mathcal L_\psi\) be the pullback of the Artin–Schreier sheaf (Lemma \hyperref[bp5]{16.1}) along the multiplication \(\mu(x,y)=xy\). Its trace function on \(\mathbf F_q\)-points is \(\psi(xy)\), and we put

\[
\operatorname{FT}_\psi K=Rp_{2!}(p_1^*K\otimes\mathcal K_\psi)[1].
\]

On trace functions it is the discrete transform up to sign: for \(y\in\mathbf F_q\),

\[
t_{\operatorname{FT}_\psi K}(y)=-\sum_{x\in\mathbf F_q}t_K(x)\,\psi(xy),
\]

because under the dictionary (Section \bpout{12.6}) pullback, tensor product and \(Rp_{2!}\) become composition, multiplication and summation over the fibres, the three operations from which \(\hat t\) is built, while the shift \([1]\) contributes the sign. Both properties of the discrete transform carry over. The first becomes the statement that the stalk at the origin of the dual line is \(R\Gamma_c(\mathbf A^1,K)[1]\) (Proposition \hyperref[bp17]{16.5}), so for \(K=j_*M\) this complex contains the group we want. Plancherel becomes an identity of trace functions (Corollary \hyperref[bp26]{16.8}), and it is what will compute the weight of the transform. In addition, Fourier inversion, \(\operatorname{FT}_{\psi^{-1}}\operatorname{FT}_\psi K\simeq K(-1)\), holds because \(R\Gamma_c(\mathbf A^1,\mathcal L_\psi)=0\), the sheaf version of \(\sum_x\psi(x)=0\). The \(\ell\)-adic transform is due to Deligne, and Laumon's proof [\hyperref[bp57]{Lau87}] of the Deligne–Laumon theorem is built on it.

To speak of the weight of the transform we need it to be a lisse sheaf on a dense open. But the transform of a sheaf need not be a sheaf. For example, the transform of the constant sheaf is a skyscraper at the origin in degree one, \(\operatorname{FT}_\psi(E)=\delta_0(-1)[-1]\). This is the sheaf version of \(\hat1=q\,\delta_0\), up to the sign coming from the shift. What the transform does preserve, however, is the category of perverse sheaves on the line (Proposition \hyperref[bp48]{16.16}, proved in Section \hyperref[bp46]{16.5} below). Since the transform is an exact equivalence of perverse sheaves, it sends the simple object \((j_*M)[1]\) to a simple object, which is either punctual or a middle extension. Moreover, only character sheaves \(\mathcal L_{\psi(ax)}\) are sent to skyscrapers. Hence, if no geometric constituent of the irreducible sheaf \(M\) is such a character, we have

\[
\operatorname{FT}_\psi((j_*M)[1])\simeq(j'_*M')[1]
\]

for an irreducible lisse \(M'\) on a dense open \(j'\) of the dual line (Proposition \hyperref[bp52]{16.17} below). Comparing stalks at the origin then identifies \(H^1_c(\mathbf A^1,j_*M)\) with a stalk of \(j'_*M'\), and this is the sheaf we asked for.

We use the six operations on the pro-étale site. At first the coefficients are a finite extension \(E/\mathbf Q_\ell\) that contains the \(p\)th roots of unity, where \(p\) is the characteristic of \(k_0\), so that \(E\) contains the values of the additive characters of \(k_0\). Later the coefficients are \(\overline{\mathbf Q}_\ell\).

% object: sec-The-Artin-Schreier-sheaf (blueprint/chap-The-Fourier-transform/sec-The-Artin-Schreier-sheaf/_section.md), content hidden
\subsection*{\texorpdfstring{16.1\quad The Artin–Schreier sheaf}{16.1 The Artin–Schreier sheaf}}
\phantomsection\label{bp3}%
\addcontentsline{toc}{subsection}{\texorpdfstring{16.1\quad The Artin–Schreier sheaf}{16.1 The Artin–Schreier sheaf}}
% text: prose-16-02
\phantomsection\label{bp4}%
The kernel of the transform is a lisse sheaf whose trace function is the additive character \(\psi\), just as \(\psi(xy)\) is the kernel of the discrete transform. It is cut out of the Artin–Schreier cover of the line as an eigenspace of the translations. This cover is Galois with group the additive group \(\mathbf F_p\), acting by \(u\mapsto u+a\). Let \(u^p-u=x\) with \(x\in\mathbf F_{q^n}\). Raising to the powers \(p,p^2,\ldots\) and adding gives \(u^{q^n}=u+\operatorname{Tr}(x)\). In other words, Frobenius acts on the fibre over \(x\) by the translation \(\operatorname{Tr}(x)\in\mathbf F_p\). A character \(\psi_0\) of \(\mathbf F_p\) then turns this translation into the number \(\psi_0(\operatorname{Tr}x)\). Similarly, the Kummer cover \(u^{q-1}=x\) of \(\mathbf G_m\) produces the multiplicative characters in the same way.

Besides the sheaf itself (Lemma \hyperref[bp5]{16.1}), we need the sheaf version of \(\sum_x\psi(x)=0\), namely the vanishing of \(R\Gamma_c(\mathbf A^r,\lambda^*\mathcal L_\psi)\) for a nonzero linear form \(\lambda\) (Lemma \hyperref[bp13]{16.4}). This is stronger than the statement on trace functions, since the trace formula only says that the alternating sum of the traces vanishes over every \(\mathbf F_{q^n}\). Fourier inversion and the stalk computations below, however, need each cohomology group to vanish. This vanishing comes from the eigenspace decomposition of the cover: the total space of the cover is again an affine line, and the translations act trivially on the compactly supported cohomology of a line. Hence the eigenspace of a nontrivial character is zero.

% object: lemma-artin-schreier-character-sheaf (blueprint/chap-The-Fourier-transform/sec-The-Artin-Schreier-sheaf/lemma-artin-schreier-character-sheaf/_section.md)
\phantomsection\label{bp5}%
\begin{bpstatementit}{Lemma~16.1}{ (The Artin–Schreier character sheaf)}
Let \(\psi:k_0\to E^\times\) be a nontrivial additive character. There is a lisse rank-one \(E\)-sheaf \(\mathcal L_\psi\) on \(\mathbf A^1_{k_0}\), with a specified trivialization at zero, such that

\[
m^*\mathcal L_\psi\simeq
\operatorname{pr}_1^*\mathcal L_\psi\otimes
\operatorname{pr}_2^*\mathcal L_\psi,
\qquad
\mathcal L_\psi^\vee\simeq\mathcal L_{\psi^{-1}}.
\]

Here \(m\) is the addition map, and the isomorphisms respect the trivializations at zero. For \(x\in\mathbf F_{q^n}\), the Frobenius trace is

\[
t_{\mathcal L_\psi,n}(x)
=\psi\bigl(\operatorname{Tr}_{\mathbf F_{q^n}/k_0}(x)\bigr).
\]

In particular \(\mathcal L_\psi\) is pure of weight zero.
\end{bpstatementit}

% proof: proof/lemma-artin-schreier-character-sheaf
\begin{proof}[{Proof}]\phantomsection\label{bp6}%
We define \(\mathcal L_\psi\) as an eigensheaf in the direct image of the constant sheaf under an Artin–Schreier cover, and read off its properties from the corresponding properties of the cover: addition of roots for the first isomorphism, a fibrewise pairing for the second, and the action of Frobenius on roots for the trace.

Fix a nontrivial character \(\psi_0:\mathbf F_p\to E^\times\). Since the trace pairing on \(k_0\) is nondegenerate, there is \(a\in k_0^\times\) such that \(\psi(z)=\psi_0(\operatorname{Tr}_{k_0/\mathbf F_p}(az))\). The cover \(u^p-u=ax\) of the affine line is finite étale, with group of translations \(\mathbf F_p\). Let \(\mathcal L_\psi\) be the image of

\[
e_{\psi_0}=\frac1p\sum_{c\in\mathbf F_p}\psi_0(c)^{-1}\tau_c^*
\]

on the direct image of the constant sheaf \(E\) under this cover. It is lisse of rank one, since its stalk is the \(\psi_0\)-eigenspace in the regular representation, and the root \(u=0\) over \(x=0\) specifies the trivialization at zero. For the first isomorphism, adding roots of two Artin–Schreier equations identifies the contracted product of the two torsors with the torsor for the sum, and passing to the character eigenspaces we obtain the isomorphism. For the second, the pairing on the direct image given by summing products over the fibres is perfect and invariant under the translations. Since, for an invariant pairing, the \(\chi\)- and \(\chi'\)-eigenspaces of two characters are orthogonal unless \(\chi\chi'=1\), our pairing pairs the \(\psi_0\)-eigenspace perfectly with the \(\psi_0^{-1}\)-eigenspace. Since \(\psi^{-1}(z)=\psi_0^{-1}(\operatorname{Tr}_{k_0/\mathbf F_p}(az))\), the \(\psi_0^{-1}\)-eigenspace of the same cover is \(\mathcal L_{\psi^{-1}}\), and we obtain \(\mathcal L_\psi^\vee\simeq\mathcal L_{\psi^{-1}}\).

For a root \(u\) above \(x\in\mathbf F_{q^n}\) we have

\[
u^{q^n}=u+
\operatorname{Tr}_{\mathbf F_{q^n}/\mathbf F_p}(ax).
\]

Geometric Frobenius acts on functions on the torsor by pullback along this map, and hence on the chosen eigenspace by \(\psi_0(\operatorname{Tr}(ax))\). By transitivity of field traces, this is the asserted trace. Since these scalars are roots of unity, their images under any field embedding have absolute value one.
\end{proof}

% text: prose-16-03
\phantomsection\label{bp7}%
We check the lemma by hand at two points. Over \(\mathbf F_p\) with \(a=1\), the fibre of the cover over \(x=1\) is \(u^p-u=1\), which has no root in \(\mathbf F_p\). Indeed, Frobenius acts on the \(p\) roots by \(u\mapsto u^p=u+1\), a translation of order \(p\), and the trace of \(\mathcal L_\psi\) at \(x=1\) is therefore \(\psi(1)=\psi_0(1)\ne1\). Over \(x=0\), by contrast, the roots are the elements of \(\mathbf F_p\), so they are all fixed and the trace is \(1\).

% object: example-gauss-sum-over-f-3 (blueprint/chap-The-Fourier-transform/sec-The-Artin-Schreier-sheaf/example-gauss-sum-over-f-3/_section.md)
\phantomsection\label{bp8}%
\begin{bpstatement}{Example~16.2}{ (A Gauss sum over $\mathbf F_3$)}
Take \(q=3\), a primitive cube root of unity \(\omega\in\overline{\mathbf Q}_\ell\) and \(\psi(x)=\omega^x\) on \(\mathbf F_3\). Pull \(\mathcal L_\psi\) back along \(h:\mathbf A^1\to\mathbf A^1\), \(x\mapsto x^2\). By the lemma, the trace of Frobenius on the stalk of \(h^*\mathcal L_\psi\) at \(x\in\mathbf F_3\) is \(\psi(x^2)\), so the trace formula gives

\[
\sum_r(-1)^r\operatorname{Tr}\bigl(F\mid H^r_c(\mathbf A^1_k,h^*\mathcal L_\psi)\bigr)
=\sum_{x\in\mathbf F_3}\omega^{x^2}=1+2\omega ,
\]

and \((1+2\omega)^2=-3\). The sheaf \(h^*\mathcal L_\psi\) is lisse of rank one. It is geometrically nontrivial, since the cover \(u^3-u=x^2\) is geometrically irreducible. Hence its geometric coinvariants vanish, and \(H^0_c=H^2_c=0\) (Lemma \bpout{15.7}). The sum is therefore \(-\operatorname{Tr}(F\mid H^1_c)\).

To see that \(H^1_c\) is a line, we use the running example. The cover \(u^3-u=x^2\) is the affine part \(Y_0\) of the elliptic curve \(C_0\colon y^2=x^3-x\) of the introduction, with the coordinates renamed (\(x^2=u^3-u\)). It lies over \(\mathbf A^1\) by \(\pi(u,x)=x\). Since \(\pi\) is the pullback along \(h\) of the Artin–Schreier cover, \(\pi_*E\) is the direct sum of \(h^*\mathcal L_{\psi'}\) over the three characters \(\psi'\) of \(\mathbf F_3\). Over \(k=\overline{\mathbf F}_3\) we obtain

\[
H^1_c(Y,E)=H^1_c(\mathbf A^1,E)\oplus H^1_c(\mathbf A^1,h^*\mathcal L_\psi)
\oplus H^1_c(\mathbf A^1,h^*\mathcal L_{\psi^{-1}}).
\]

The left side is \(H^1(C,E)\), because \(C\) minus \(Y\) is the single point at infinity, and it has dimension \(2\). Since the first summand on the right is zero, it remains to see that the last two summands have the same dimension. For this, note that the automorphism \((u,x)\mapsto(-u,ix)\) of \(Y\), with \(i^2=-1\) in \(\mathbf F_9\), conjugates the translation by \(c\) into the translation by \(-c\). It therefore exchanges the last two summands, and each of them is a line. This automorphism is defined only over \(\mathbf F_9\), but we use it only to count dimensions. The decomposition itself is stable under Frobenius, because the translations are defined over \(\mathbf F_3\).

Thus Frobenius acts on the line \(H^1_c(\mathbf A^1,h^*\mathcal L_\psi)\) by \(-(1+2\omega)\). This is a square root of \(-3\), and it is one of the two reciprocal roots \(\pm i\sqrt3\) of \(P_1(T)=1+3T^2\) found in the introduction. Moreover, the number \(1+2\omega=\omega-\omega^2\) is the Gauss sum \(g\) of the quadratic character of \(\mathbf F_3\). The classical fact \(|g|=\sqrt q\) is thus the statement that this \(H^1_c\) is pure of weight one. Chapter \bpout{17} proves the analogue for every such sheaf.
\end{bpstatement}

% text: prose-16-04
\phantomsection\label{bp9}%
Fourier inversion integrates the kernel over a whole fibre \(\mathbf A^r\) of a vector bundle, so we need the vanishing on \(\mathbf A^r\) and not only on the line. In coordinates in which the linear form is the first coordinate, the integral over \(\mathbf A^r\) is a product of an integral over the line and one over \(\mathbf A^{r-1}\), and the Künneth formula makes this precise.

Let \(f:X\to S\) and \(g:Y\to S\) be separated morphisms of finite type, let \(p_1:X\times_SY\to X\) and \(p_2:X\times_SY\to Y\) be the projections, and let \(h=f\circ p_1:X\times_SY\to S\) be the structure map. For complexes \(K\) on \(X\) and \(L\) on \(Y\) we put \(K\boxtimes_SL=p_1^*K\otimes p_2^*L\).

% object: lemma-kunneth-formula-compact-supports (blueprint/chap-The-Fourier-transform/sec-The-Artin-Schreier-sheaf/lemma-kunneth-formula-compact-supports/_section.md)
\phantomsection\label{bp10}%
\begin{bpstatementit}{Lemma~16.3}{ (Künneth formula for compact supports)}
In this situation, let \(K,L\) be bounded constructible complexes with coefficients in a field. Then the maps in

\[
Rh_!(p_1^*K\otimes p_2^*L)
\simeq Rf_!(K\otimes Rp_{1!}p_2^*L)
\simeq Rf_!(K\otimes f^*Rg_!L)
\simeq Rf_!K\otimes Rg_!L
\]

of the projection formula for \(p_1\), the proper base change isomorphism \(Rp_{1!}p_2^*L\to f^*Rg_!L\) and the projection formula for \(f\) are isomorphisms. We call their composite \(Rh_!(K\boxtimes_SL)\simeq Rf_!K\otimes Rg_!L\) the \emph{Künneth isomorphism}.
\end{bpstatementit}

% proof: proof/lemma-kunneth-formula-compact-supports
\begin{proof}[{Proof}]\phantomsection\label{bp11}%
We integrate first along the fibres of \(p_1\), which are copies of the fibres of \(g\). This is possible because \(h=f\circ p_1\), so that \(Rh_!=Rf_!Rp_{1!}\). Since the square formed by \(p_1\), \(p_2\), \(f\) and \(g\) is cartesian, proper base change gives \(Rp_{1!}p_2^*L\simeq f^*Rg_!L\). The two projection formulas are isomorphisms as well, and hence so is the composite.
\end{proof}

% text: prose-16-05
\phantomsection\label{bp12}%
Over a point, this is the Künneth isomorphism for cohomology with compact support. All three isomorphisms in the composite are natural, so when \(X\), \(Y\), \(S\), \(f\), \(g\), \(K\) and \(L\) are defined over \(k_0\), the Künneth isomorphism commutes with the Weil structures.

The case \(\lambda=0\) of the next lemma is needed as well, since Fourier inversion integrates the kernel along the diagonal, where the character is trivial.

% object: lemma-integral-linear-character (blueprint/chap-The-Fourier-transform/sec-The-Artin-Schreier-sheaf/lemma-integral-linear-character/_section.md)
\phantomsection\label{bp13}%
\begin{bpstatementit}{Lemma~16.4}{ (The integral of a linear character)}
Let \(\lambda\) be a linear form on \(\mathbf A^r\) over an algebraically closed field of characteristic \(p\). Then

\[
R\Gamma_c(\mathbf A^r,\lambda^*\mathcal L_\psi)
\simeq
\begin{cases}
E(-r)[-2r],&\lambda=0,\\
0,&\lambda\ne0.
\end{cases}
\]

For \(\lambda=0\), the isomorphism is the trace map for cohomology with compact support.
\end{bpstatementit}

% proof: proof/lemma-integral-linear-character
\begin{proof}[{Proof}]\phantomsection\label{bp14}%
On the line, from the cohomology of \(\mathbf P^1\) and localization at the point at infinity, we get \(R\Gamma_c(\mathbf A^1,E)=E(-1)[-2]\), identified by the trace map (Appendix \bpout{D}). Every translation acts trivially on it, since it preserves the trace and the cohomology is concentrated in degree two. Now let \(\pi:\mathbf A^1\to\mathbf A^1\), \(u\mapsto u^p-u\), be the Artin–Schreier cover, which is finite étale and whose total space is again an affine line. Let \(\mathcal L_{\psi_0}\) be the \(\psi_0\)-eigenspace of \(\pi_*E\), that is, the sheaf of the proof of Lemma \hyperref[bp5]{16.1} for \(a=1\). The projector \(e_{\psi_0}\) of that proof cuts out \(R\Gamma_c(\mathbf A^1,\mathcal L_{\psi_0})\) from \(R\Gamma_c(\mathbf A^1,\pi_*E)=R\Gamma_c(\mathbf A^1,E)\) as the \(\psi_0\)-eigenspace of the deck translations. As we have just seen, these translations act trivially on the cohomology of the upper line, so, since \(\psi_0\) is nontrivial, the eigenspace is zero, that is, \(R\Gamma_c(\mathbf A^1,\mathcal L_{\psi_0})=0\). Finally, by its construction, \(\mathcal L_\psi\) is the pullback of \(\mathcal L_{\psi_0}\) by the automorphism \(x\mapsto ax\) of the line, and hence \(R\Gamma_c(\mathbf A^1,\mathcal L_\psi)=0\) as well.

On \(\mathbf A^r\) with \(\lambda\ne0\), we choose linear coordinates in which \(\lambda\) is the first coordinate. Then \(\lambda^*\mathcal L_\psi\) is the external product of \(\mathcal L_\psi\) on the first factor with the constant sheaf on \(\mathbf A^{r-1}\), and the Künneth isomorphism of Lemma \hyperref[bp10]{16.3} writes its cohomology as a tensor product with the first factor \(R\Gamma_c(\mathbf A^1,\mathcal L_\psi)=0\). For \(\lambda=0\), on the other hand, the sheaf is constant, so the trace map of Definition \bpout{9.24} for the projection \(\mathbf A^r\to\operatorname{Spec}k\) is an isomorphism \(R\Gamma_c(\mathbf A^r,E)(r)[2r]\simeq E\). On constant coefficients this map is the trace map of affine space (9.11), obtained by integrating along one coordinate line at a time. Finally, all these calculations take place with coefficients in \(E\) and commute with extension of scalars.
\end{proof}


% object: sec-Fourier-inversion (blueprint/chap-The-Fourier-transform/sec-Fourier-inversion/_section.md), content hidden
\subsection*{\texorpdfstring{16.2\quad Fourier inversion}{16.2 Fourier inversion}}
\phantomsection\label{bp15}%
\addcontentsline{toc}{subsection}{\texorpdfstring{16.2\quad Fourier inversion}{16.2 Fourier inversion}}
% text: prose-16-06
\phantomsection\label{bp16}%
We integrate with compact supports, \(Rp_{2!}\) rather than \(Rp_{2*}\), for two reasons. First, proper base change computes the stalks of \(Rp_{2!}\) as compactly supported cohomology of the fibres, and this puts \(H^1_c(\mathbf A^1,j_*M)\) at the origin. Second, the trace formula holds for \(Rp_{2!}\), so that on trace functions the transform becomes the discrete Fourier transform (16.3) (proved in Section \hyperref[bp24]{16.3}).

The shift \([r]\) is the one for which the transform preserves the perverse heart. For example, on the line we will find \(\operatorname{FT}_\psi(E[1])=\delta_0(-1)\), whereas without the shift the perverse sheaf \(E[1]\) would go to \(\delta_0(-1)[-1]\), which is not perverse.

We work with a vector bundle \(V\to S\) rather than the line alone, because the proofs move along the base. For instance, the stalk formula is base change to a point of the dual bundle, and Fourier inversion is a computation on \(V\times_SV^\vee\times_SV\), which reduces, fibre by fibre, to Lemma \hyperref[bp13]{16.4} on \(\mathbf A^r\). We prove these two properties in Propositions \hyperref[bp17]{16.5} and \hyperref[bp19]{Proposition 16.6}, while the preservation of the heart is Proposition \hyperref[bp48]{16.16}.

Let \(S\) be a scheme of finite type over \(k_0\) or over \(k\), let \(V\to S\) be a vector bundle of rank \(r\), and let \(V^\vee\) be its dual bundle. Write \(p_1,p_2\) for the projections of \(V\times_SV^\vee\), and set

\[
\mathcal K_\psi=\langle\ ,\ \rangle^*\mathcal L_\psi,
\qquad
\operatorname{FT}_\psi(K)
=Rp_{2!}(p_1^*K\otimes\mathcal K_\psi)[r].
\tag{16.1}
\]

The kernel \(\mathcal K_\psi\) is lisse of rank one, and since the six operations preserve constructibility, this defines a triangulated functor \(\operatorname{FT}_\psi:\Dbc(V,E)\to\Dbc(V^\vee,E)\), normalized by the shift \([r]\).

% object: proposition-base-change-fourier-inversion (blueprint/chap-The-Fourier-transform/sec-Fourier-inversion/proposition-base-change-fourier-inversion/_section.md)
\phantomsection\label{bp17}%
\begin{bpstatementit}{Proposition~16.5}{ (Base change and stalks)}
Let \(g:S'\to S\) be a morphism of schemes with \(S'\) of finite type over \(k_0\) or over \(k\), let \(V'=V\times_SS'\), and let \(g_V:V'\to V\) and \(g_{V^\vee}:V'^\vee\to V^\vee\) be the projections.

\begin{enumerate}
\item The proper base change isomorphism for \(p_2\) induces an isomorphism \(g_{V^\vee}^*\operatorname{FT}_\psi(K)\simeq
  \operatorname{FT}_\psi(g_V^*K)\).
\item For \(N\in\Dbc(S,E)\), the projection formula for \(p_2\) induces an isomorphism \(\operatorname{FT}_\psi(K\otimes\pi^*N)\simeq
  \operatorname{FT}_\psi(K)\otimes\pi^{\vee*}N\), where \(\pi:V\to S\) and \(\pi^\vee:V^\vee\to S\) are the projections.
\item For a geometric point \(\bar y\) of \(V^\vee\), the proper base change map gives \[
   (\operatorname{FT}_\psi K)_{\bar y}
   \simeq
   R\Gamma_c\bigl(V_{\bar y},
   K|_{V_{\bar y}}\otimes
   \mathcal L_{\psi\langle x,y\rangle}\bigr)[r].
   \tag{16.2}
   \]
\end{enumerate}
\end{bpstatementit}

% proof: proof/proposition-base-change-fourier-inversion
\begin{proof}[{Proof}]\phantomsection\label{bp18}%
For (1), the formation of \(V\times_SV^\vee\), of the pairing and of \(\mathcal K_\psi\) commutes with base change along \(g\), and proper base change for \(Rp_{2!}\) gives the isomorphism. For (2), we have \(p_1^*\pi^*N=p_2^*\pi^{\vee*}N\), and the projection formula for \(p_2\) gives the isomorphism. For (3), the fibre of \(p_2\) over \(\bar y\) is \(V_{\bar y}\), on which \(\mathcal K_\psi\) restricts to \(\mathcal L_{\psi\langle x,y\rangle}\), and proper base change for \(Rp_{2!}\) along \(\bar y\to V^\vee\) gives the formula.
\end{proof}

% object: proposition-fourier-inversion (blueprint/chap-The-Fourier-transform/sec-Fourier-inversion/proposition-fourier-inversion/_section.md)
\phantomsection\label{bp19}%
\begin{bpstatementit}{Proposition~16.6}{ (Fourier inversion)}
Under \(V^{\vee\vee}=V\) there are natural isomorphisms of functors \(\Dbc(V,E)\to\Dbc(V,E)\)

\[
\operatorname{FT}_{\psi^{-1}}\operatorname{FT}_\psi(K)
\simeq K(-r),
\qquad
\operatorname{FT}_{\psi}\operatorname{FT}_\psi(K)
\simeq [-1]_V^*K(-r).
\]
\end{bpstatementit}

% proof: proof/proposition-fourier-inversion
\begin{proof}[{Proof}]\phantomsection\label{bp20}%
The composite of two transforms is again an integral transform, and we compute its kernel. It will turn out to be the constant sheaf on the diagonal, twisted by \((-r)\), which up to this twist is the kernel of the identity functor. Let \(r_{ij}\) be the projections of \(V\times_SV^\vee\times_SV\) onto pairs of factors. By base change and the projection formula, the composite of two transforms is the transform with the convolution kernel

\[
Rr_{13!}(r_{12}^*\mathcal K_\psi\otimes
r_{23}^*\mathcal K_{\psi^{-1}})[2r].
\]

At a point \((x,z)\), the unshifted integral over the middle factor is the cohomology of the character sheaf attached to \(y\mapsto\psi\langle x-z,y\rangle\), and hence it vanishes for \(x\ne z\) by Lemma \hyperref[bp13]{16.4}. Thus the unshifted convolution is supported on the diagonal, so the localization triangle identifies it canonically with the direct image of its restriction to the diagonal. There the unit trivialization of the character sheaf and the trace map for the vector bundle identify the restriction with \(E(-r)[-2r]\). Since the shift \([2r]\) cancels \([-2r]\), we obtain the first inversion formula. If instead the two kernels use the same character, then \(x-z\) is replaced by \(x+z\), the support is the antidiagonal, and we obtain the second formula. Since the traces and unit trivializations used here are defined over \(k_0\), these isomorphisms commute with the Weil structures.
\end{proof}

% text: prose-16-07
\phantomsection\label{bp21}%
On the line, the transform exchanges the two simplest perverse sheaves, the shifted constant sheaf and the skyscraper at the origin, up to a twist. This is the sheaf version of \(\hat 1=q\,\delta_0\) and \(\hat\delta_0=1\).

% object: lemma-fourier-transform-constant-sheaf (blueprint/chap-The-Fourier-transform/sec-Fourier-inversion/lemma-fourier-transform-constant-sheaf/_section.md)
\phantomsection\label{bp22}%
\begin{bpstatementit}{Lemma~16.7}{ (Transform of the constant sheaf)}
Let \(V=\mathbf A^1_{k_0}\) over \(S=\operatorname{Spec}k_0\), and let \(\delta_0\) be the constant skyscraper sheaf \(E\) at the origin. There are isomorphisms of Weil complexes

\[
\operatorname{FT}_\psi(E)\simeq\delta_0(-1)[-1],
\qquad \operatorname{FT}_\psi(\delta_0)\simeq E[1].
\]
\end{bpstatementit}

% proof: proof/lemma-fourier-transform-constant-sheaf
\begin{proof}[{Proof}]\phantomsection\label{bp23}%
For \(K=E\), the stalk formula (16.2) at \(y\) is the integral of the character sheaf \(\mathcal L_{\psi(xy)}\) over the line, shifted by one. By Lemma \hyperref[bp13]{16.4} it vanishes for \(y\ne0\) and equals \(E(-1)[-2]\) at \(y=0\). Hence \(\operatorname{FT}_\psi(E)\) is supported at the origin, and with the shift \([1]\) it is \(\delta_0(-1)[-1]\). For \(K=\delta_0\), the pullback by the first projection is supported on \(\{0\}\times\mathbf A^1\). Since the kernel is trivial there and the second projection is an isomorphism, this gives \(E[1]\). Moreover, the trace map and the unit trivialization are defined over \(k_0\), so both identifications are Frobenius-equivariant.
\end{proof}


% object: sec-The-Plancherel-formula (blueprint/chap-The-Fourier-transform/sec-The-Plancherel-formula/_section.md), content hidden
\subsection*{\texorpdfstring{16.3\quad The Plancherel formula}{16.3 The Plancherel formula}}
\phantomsection\label{bp24}%
\addcontentsline{toc}{subsection}{\texorpdfstring{16.3\quad The Plancherel formula}{16.3 The Plancherel formula}}
% text: prose-16-08
\phantomsection\label{bp25}%
On trace functions, the Fourier transform is, up to the sign \((-1)^r\), the discrete Fourier transform over \(\mathbf F_{q^n}\) for the character \(\psi\circ\operatorname{Tr}_{\mathbf F_{q^n}/k_0}\) (equation (16.3) below). The Plancherel formula then says how it changes the growth of squared trace sums, and for a pure sheaf this growth is the weight (Theorem \bpout{15.14}). This is how we find the weight of the transform in Chapter \bpout{17}: on the line, Plancherel multiplies the sums of squared traces by \(q^n\), and so it raises the weight read off from their growth by one.

In this section \(S=\operatorname{Spec}k_0\) and \(V=\mathbf A^r_{k_0}\). For a Weil complex \(K\) on a \(k_0\)-scheme \(X_0\), we write \(t_{K,n}(x)\) for the alternating sum of the Frobenius traces at \(x\in X_0(\mathbf F_{q^n})\), as in Definition \bpout{15.10}. With this notation, the dictionary (Proposition \bpout{12.28}) computes the trace function of the transform step by step. For \(y\in\mathbf F_{q^n}^r\), the shift contributes \((-1)^r\) (chain formula of Proposition \bpout{11.2}), part (3) sums over the fibre of \(p_2\), and part (2) splits the trace of the tensor product:

\[
t_{\operatorname{FT}_\psi K,n}(y)
=(-1)^r\,t_{Rp_{2!}(p_1^*K\otimes\mathcal K_\psi),n}(y)
=(-1)^r\sum_{x\in\mathbf F_{q^n}^r}t_{p_1^*K,n}(x,y)\,
t_{\mathcal K_\psi,n}(x,y).
\]

By part (1), \(t_{p_1^*K,n}(x,y)=t_{K,n}(x)\), and by part (1) and Lemma \hyperref[bp5]{16.1}, \(t_{\mathcal K_\psi,n}(x,y)=\psi(\operatorname{Tr}_{\mathbf F_{q^n}/k_0}\langle x,y\rangle)\). Hence

\[
t_{\operatorname{FT}_\psi K,n}(y)
=(-1)^r\sum_{x\in\mathbf F_{q^n}^r}
t_{K,n}(x)\,
\psi\bigl(\operatorname{Tr}_{\mathbf F_{q^n}/k_0}
\langle x,y\rangle\bigr).
\tag{16.3}
\]

% object: corollary-plancherel (blueprint/chap-The-Fourier-transform/sec-The-Plancherel-formula/corollary-plancherel/_section.md)
\phantomsection\label{bp26}%
\begin{bpstatementit}{Corollary~16.8}{ (Plancherel)}
For every \(n\ge1\),

\[
\sum_y|\iota t_{\operatorname{FT}_\psi K,n}(y)|^2
=q^{nr}\sum_x|\iota t_{K,n}(x)|^2.
\]

Consequently, for \(r=1\), the number \(\nu\) of Definition \bpout{15.10} satisfies \(\nu(\operatorname{FT}_\psi K)=\nu(K)+1\) whenever these numbers are finite.
\end{bpstatementit}

% proof: proof/corollary-plancherel
\begin{proof}[{Proof}]\phantomsection\label{bp27}%
As for the classical transform, the proof is by orthogonality of characters. We expand the square of the sum in (16.3), using that the complex conjugate of a character value is its inverse. This writes the left side as a double sum over \(x,z\in\mathbf F_{q^n}^r\) of \(\iota t_{K,n}(x)\overline{\iota t_{K,n}(z)}\) times a character sum over \(y\), which we now compute. For \(x,z\in\mathbf F_{q^n}^r\), the orthogonality of characters gives

\[
\sum_y\iota\psi\bigl(\operatorname{Tr}\langle x-z,y\rangle\bigr)
=\begin{cases}q^{nr},&x=z,\\0,&x\ne z.\end{cases}
\]

Indeed, for \(x\ne z\) the linear form \(y\mapsto\operatorname{Tr}_{\mathbf F_{q^n}/k_0}\langle x-z,y\rangle\) is nonzero, since the trace pairing is nondegenerate, and hence surjective onto \(k_0\). As \(\psi\) is nontrivial, there is \(y_0\) with \(c=\psi(\operatorname{Tr}\langle x-z,y_0\rangle)\ne1\), and substituting \(y+y_0\) for \(y\) multiplies the sum by \(\iota(c)\), so the sum is zero. Therefore only the diagonal terms remain, and they give the asserted identity. On a line, finally, the factor \(q^n\) raises \(\nu\) by one.
\end{proof}

% text: prose-16-09
\phantomsection\label{bp28}%
For the constant sheaf \(E\) on the line and \(n=1\), formula (16.3) gives the trace \(-\sum_x\psi(xy)\), which is \(-q\) at \(y=0\) and \(0\) elsewhere, in agreement with \(\operatorname{FT}_\psi(E)=\delta_0(-1)[-1]\). Both sides of Plancherel are then \(q^2\), and over \(\mathbf F_{q^n}\) they are \(q^{2n}\). This gives \(\nu(E)=0\) and \(\nu(\operatorname{FT}_\psi E)=1\), as the corollary predicts.


% object: sec-Perverse-sheaves-on-a-curve (blueprint/chap-The-Fourier-transform/sec-Perverse-sheaves-on-a-curve/_section.md), content hidden
\subsection*{\texorpdfstring{16.4\quad Perverse sheaves on a curve}{16.4 Perverse sheaves on a curve}}
\phantomsection\label{bp29}%
\addcontentsline{toc}{subsection}{\texorpdfstring{16.4\quad Perverse sheaves on a curve}{16.4 Perverse sheaves on a curve}}
% text: prose-16-10
\phantomsection\label{bp30}%
To apply Theorem \bpout{15.14} to the transform of a pure lisse sheaf \(M\), we need to know that the transform of \((j_*M)[1]\) is again a lisse sheaf on a dense open, placed in degree \(-1\). Although the transform sends the constant sheaf to a skyscraper (Lemma \hyperref[bp22]{16.7}), it preserves the perverse heart, which on a smooth curve has an elementary description in terms of ordinary cohomology sheaves and punctual support.

We can read off the shape of this description from the stalk formula. The stalk of the transform of \(G[1]\) at \(y\) is \(R\Gamma_c(\mathbf A^1,G\otimes\mathcal L_{\psi(yx)})[2]\). For a sheaf \(G\) without punctual sections \(H^0_c\) vanishes, and the proof of Proposition \hyperref[bp48]{16.16} below shows that \(H^2_c\) is nonzero only at finitely many \(y\). So the transform lives in degrees \(-1\) and \(0\), with a punctual sheaf in degree \(0\), and this suggests allowing skyscrapers in degree \(0\) in \({}^pD^{\le0}\). Orthogonality then forces the condition on the other side: a sheaf in degree \(0\) lies in \({}^pD^{\ge1}\) only if it receives no nonzero map from a skyscraper, that is, only if it has no punctual subsheaf. The perverse t-structure on a curve is the t-structure built from these two conditions.

What we finally use about it is the list of its simple objects (Lemma \hyperref[bp43]{16.14}). By this list, an exact equivalence must send an irreducible middle extension to an irreducible middle extension or to an irreducible punctual sheaf.

Recall that a \emph{t-structure} on a triangulated category \(D\) is a pair of full subcategories \(D^{\le0}\supset D^{\le-1}\) and \(D^{\ge0}\supset D^{\ge1}\), where \(D^{\le n}=D^{\le0}[-n]\) and \(D^{\ge n}=D^{\ge0}[-n]\), such that \(\Hom(A,B)=0\) for \(A\in D^{\le0}\) and \(B\in D^{\ge1}\), and every object sits in a triangle \(A\to K\to
B\to\) with \(A\in D^{\le0}\) and \(B\in D^{\ge1}\). Its \emph{heart} is \(D^{\le0}\cap D^{\ge0}\); for the ordinary t-structure it is the category of sheaves. We write \(D^{[a,b]}=D^{\ge a}\cap D^{\le b}\). The triangle is unique up to unique isomorphism, so, as for the ordinary t-structure, it gives truncation functors and cohomology objects of \(K\) in the heart. Finally, a functor is \emph{t-exact} if it preserves both subcategories.

Let \(C\) be a smooth curve over a perfect field. On \(C\) the perverse t-structure is built from punctual sheaves.

% object: definition-punctual-sheaf (blueprint/chap-The-Fourier-transform/sec-Perverse-sheaves-on-a-curve/definition-punctual-sheaf/_section.md)
\phantomsection\label{bp31}%
\begin{bpstatement}{Definition~16.9}{ (Punctual sheaf)}
A constructible sheaf on \(C\) is \emph{punctual} if its support is a finite set of closed points.
\end{bpstatement}

% text: prose-16-11
\phantomsection\label{bp32}%
The truncation triangle of the perverse t-structure splits a sheaf into a punctual part and a part without punctual subsheaves, so we first construct the largest punctual subsheaf. In the following lemma “sheaf” means a constructible sheaf with coefficients \(E\) or \(\overline{\mathbf Q}_\ell\).

% object: lemma-largest-punctual-subsheaf (blueprint/chap-The-Fourier-transform/sec-Perverse-sheaves-on-a-curve/lemma-largest-punctual-subsheaf/_section.md)
\phantomsection\label{bp33}%
\begin{bpstatementit}{Lemma~16.10}{ (The largest punctual subsheaf)}
A lisse sheaf on \(C\) has no nonzero punctual subsheaf. Let \(F\) be a constructible sheaf on \(C\), let \(j:U\hookrightarrow C\) be a dense open on which \(F\) is lisse, and put

\[
F_{\rm p}=\ker(F\longrightarrow j_*j^*F).
\]

Then \(F_{\rm p}\) is the largest punctual subsheaf of \(F\), and \(F/F_{\rm p}\) has no nonzero punctual subsheaf. In particular, \(F_{\rm p}\) does not depend on \(U\).
\end{bpstatementit}

% proof: proof/lemma-largest-punctual-subsheaf
\begin{proof}[{Proof}]\phantomsection\label{bp34}%
Let \(L\) be a lisse sheaf. For a geometric point \(\bar x\) and a geometric generic point \(\bar\eta\) specializing to it, the generization map \(L_{\bar x}\to L_{\bar\eta}\) is injective. Indeed, for a sheaf defined over a finite extension \(E\) of \(\mathbf Q_\ell\), this holds for the finite locally constant reductions of the lattice in Proposition \bpout{13.10}, and we pass to the inverse limit, invert the uniformizer and extend scalars. Now a punctual subsheaf \(G\subset L\) has \(G_{\bar\eta}=0\), and since the generization maps of \(G\) and \(L\) are compatible, \(G_{\bar x}=0\).

Now let \(F\) and \(U\) be as in the statement. The sheaf \(F_{\rm p}\) is supported on the finite set \(C\setminus U\), hence punctual. Conversely, a punctual subsheaf \(G\) of \(F\) restricts to a punctual subsheaf of the lisse sheaf \(j^*F\), hence to zero. Since its map to \(j_*j^*F\) is adjoint to \(j^*G\to j^*F\), this map vanishes, and \(G\subset F_{\rm p}\). Finally, \(j_*j^*F\) has no nonzero punctual subsheaf. Indeed, if \(H\subset j_*j^*F\) is punctual, then \(j^*H=0\), and the inclusion is zero since \(\Hom(H,j_*j^*F)=\Hom(j^*H,j^*F)=0\). As \(F/F_{\rm p}\) embeds in \(j_*j^*F\), it has no nonzero punctual subsheaf either.
\end{proof}

% object: definition-perverse-sheaves-curve (blueprint/chap-The-Fourier-transform/sec-Perverse-sheaves-on-a-curve/definition-perverse-sheaves-curve/_section.md)
\phantomsection\label{bp35}%
\begin{bpstatement}{Definition~16.11}{ (Perverse sheaves on a curve)}
Let \(\Lambda\) be \(E\) or \(\overline{\mathbf Q}_\ell\). Define two classes of objects of \(D^b_c(C_{\mathrm{pro}\acute{e}t},\Lambda)\) by

\[
\begin{split}
{}^pD^{\le0}&=\{K:\mathcal H^aK=0\ (a>0),\
\mathcal H^0K\text{ is punctual}\},\\
{}^pD^{\ge0}&=\{K:\mathcal H^aK=0\ (a<-1),\
\mathcal H^{-1}K\text{ has no nonzero punctual subsheaf}\},
\end{split}
\]

and \({}^pD^{\le n}={}^pD^{\le0}[-n]\), \({}^pD^{\ge n}={}^pD^{\ge0}[-n]\). A \emph{perverse sheaf} on \(C\) is an object of \({}^pD^{\le0}\cap{}^pD^{\ge0}\). For \(C=C_0\times_{k_0}k\), a \emph{perverse Weil sheaf} on \(C_0\) is a perverse sheaf on \(C\) with a Weil structure, with morphisms as in Definition \bpout{10.18}.
\end{bpstatement}

% text: prose-16-12
\phantomsection\label{bp36}%
By Lemma \hyperref[bp37]{16.12} below these classes form a t-structure, the perverse t-structure. We write \({}^pD^{[a,b]}\) and \({}^p\mathcal H^n\) for its subcategories and cohomology objects.

On the line, the two basic examples are \(E[1]\) and the skyscraper \(\delta_0\), and by Lemma \hyperref[bp22]{16.7} the Fourier transform exchanges them up to a Tate twist. This is the simplest instance of the exactness proved in Proposition \hyperref[bp48]{16.16}.

% object: lemma-perverse-sheaves-curve (blueprint/chap-The-Fourier-transform/sec-Perverse-sheaves-on-a-curve/lemma-perverse-sheaves-curve/_section.md)
\phantomsection\label{bp37}%
\begin{bpstatementit}{Lemma~16.12}{ (Perverse sheaves on a curve)}
The classes \({}^pD^{\le0}\) and \({}^pD^{\ge0}\) of Definition \hyperref[bp35]{16.11} form a bounded t-structure on \(D^b_c(C_{\mathrm{pro}\acute{e}t},\overline{\mathbf Q}_\ell)\), and its heart, the category of perverse sheaves on \(C\), is abelian. A punctual sheaf in degree zero is perverse, and so is \((j_*M)[1]\) for a lisse sheaf \(M\) on a dense open \(j:U\hookrightarrow C\).
\end{bpstatementit}

% proof: proof/lemma-perverse-sheaves-curve
\begin{proof}[{Proof}]\phantomsection\label{bp38}%
The truncation triangle of \(K\) splits \(\mathcal H^0K\) into its largest punctual subsheaf, which goes into \({}^pD^{\le0}\), and the quotient, which goes into \({}^pD^{\ge1}\), using Lemma \hyperref[bp33]{16.10}.

The shift axioms are immediate from the definitions. For orthogonality, if \(A\in{}^pD^{\le0}\) and \(B\in{}^pD^{\ge1}\), then ordinary truncation identifies \(\Hom(A,B)\) with \(\Hom(\mathcal H^0A,\mathcal H^0B)\), which is zero because the source is punctual and the target has no punctual subsheaf. To construct the truncation triangle for \(K\), put \(K'=\tau^{\le0}K\) and \(F=\mathcal H^0K\), and let \(A\) be the fibre of \(K'\to F/F_{\rm p}\). Then the cohomology of \(A\) is that of \(K\) below zero, and \(\mathcal H^0A=F_{\rm p}\), so that \(A\in{}^pD^{\le0}\). By the octahedral axiom, applied to \(A\to K'\to K\), we obtain a triangle \(A\to K\to B\to\), in which the cone \(B\) is an extension of \(\tau^{\ge1}K\) by \(F/F_{\rm p}\) in degree zero, so it has no negative ordinary cohomology and \(\mathcal H^0B=F/F_{\rm p}\). Since \(F/F_{\rm p}\) has no punctual subsheaf by Lemma \hyperref[bp33]{16.10}, \(B\in{}^pD^{\ge1}\). Finally, the t-structure is bounded because the ordinary one is.

The heart of any t-structure is abelian, with kernels and cokernels computed from the perverse cohomology of a cone [\hyperref[bp58]{BBD82}, Théorème 1.3.6]; the proof uses only the octahedral axiom and orthogonality. Punctual sheaves in degree zero clearly lie in the heart. So does \((j_*M)[1]\), since \(j_*M\) has no punctual subsheaf, as follows from Lemma \hyperref[bp33]{16.10} applied to \(F=j_*M\), for which \(F_{\rm p}=0\).
\end{proof}

% text: prose-16-13
\phantomsection\label{bp39}%
Unwinding the definition, a perverse sheaf on \(C\) is a complex in ordinary degrees \(-1,0\) whose \(\mathcal H^0\) is punctual and whose \(\mathcal H^{-1}\) has no nonzero punctual subsheaf.

% object: lemma-perverse-weil-sheaves (blueprint/chap-The-Fourier-transform/sec-Perverse-sheaves-on-a-curve/lemma-perverse-weil-sheaves/_section.md)
\phantomsection\label{bp40}%
\begin{bpstatementit}{Lemma~16.13}{ (Perverse Weil sheaves)}
The perverse Weil sheaves on \(C_0\) form an abelian category, and the functor forgetting the Weil structure, to perverse sheaves on \(C\), is exact and reflects exactness.
\end{bpstatementit}

% proof: proof/lemma-perverse-weil-sheaves
\begin{proof}[{Proof}]\phantomsection\label{bp41}%
Frobenius pullback is an exact equivalence for the ordinary t-structure, and it preserves punctual support and the absence of punctual subsheaves. It is therefore t-exact for the perverse t-structure, which means that it carries the kernel and cokernel of a compatible morphism of perverse sheaves to the kernel and cokernel of the pulled-back morphism. The Weil isomorphisms and their inverses therefore induce inverse isomorphisms on the kernel and cokernel, and by naturality the coimage–image map is compatible with them. So perverse Weil sheaves form an abelian category, with kernels and cokernels computed in perverse sheaves on \(C\), and this is the assertion about the forgetful functor.
\end{proof}

% text: prose-16-14
\phantomsection\label{bp42}%
An exact equivalence of abelian categories sends simple objects to simple objects. Once we know that the transform is exact on the heart, the following list of simple objects will tell us what the transform of a middle extension can be.

% object: lemma-simple-objects (blueprint/chap-The-Fourier-transform/sec-Perverse-sheaves-on-a-curve/lemma-simple-objects/_section.md)
\phantomsection\label{bp43}%
\begin{bpstatementit}{Lemma~16.14}{ (The simple objects)}
In the category of perverse sheaves on \(C\), and in that of perverse Weil sheaves on \(C_0\), the simple objects are the irreducible punctual objects and the objects \((j_*M)[1]\) with \(M\) nonzero irreducible lisse on a dense open, where irreducibility is taken in the geometric or in the Weil category, respectively. These middle extensions have no nonzero punctual subobject or quotient in the perverse heart.
\end{bpstatementit}

% proof: proof/lemma-simple-objects
\begin{proof}[{Proof}]\phantomsection\label{bp44}%
Since a middle extension has no punctual subobject or quotient and \(M\) is irreducible on the open, the middle extension is simple. Conversely, ordinary truncation splits an arbitrary perverse sheaf into a punctual part and a part of the form \(F[1]\), and simplicity forces one of them to vanish.

We first show that \(P=(j_*M)[1]\) has no punctual subobject or quotient. The adjunction calculations below are geometric, so they apply equally to compatible Weil morphisms. Let \(i\) be the inclusion of a finite set of closed points and \(A\) a sheaf on it. A map \(P\to i_*A\) is zero, because by adjunction it corresponds to a map from \(i^*P=(i^*j_*M)[1]\) to the sheaf \(A\), that is, to a negative Ext. A map in the other direction is also zero, and to see this we first enlarge the boundary. Since \(M\) is lisse, \(j_*M\) does not change when we replace \(U\) by a smaller dense open and \(M\) by its restriction, so we may assume that the support of \(A\) lies in \(Z=C\setminus U\), and we let \(i\) now denote the inclusion of \(Z\). Then the adjunction \(j^*\dashv Rj_*\) shows that \(i^!Rj_*M=0\). In the triangle

\[
j_*M\longrightarrow Rj_*M\longrightarrow
\tau^{\ge1}Rj_*M\longrightarrow
\]

the last term is supported on \(Z\), so that \(i^!\) agrees with \(i^*\) on it. Applying \(i^!\) and using \(i^!Rj_*M=0\), we find \(i^!j_*M\simeq(i^*\tau^{\ge1}Rj_*M)[-1]\), which has ordinary cohomology only in degrees at least two. Hence \(i^!P\) has ordinary cohomology only in degrees at least one. Therefore \(\Hom(i_*A,P)=\Hom(A,i^!P)=0\).

To see that \(P\) is simple, note that on a dense open a perverse subobject of \(P\) is a shifted lisse subsheaf of \(M\). By irreducibility it is \(0\) or everything there, and the difference is then punctual. In detail, consider a perverse exact sequence \(0\to Q\to P\to R\to0\), restrict it to a smaller dense open on which all ordinary cohomology sheaves are lisse. There \(P=M[1]\), and the sequence is a shifted exact sequence of lisse sheaves. Moreover, restriction to a dense open preserves irreducibility, because the fundamental group of a dense open of a normal connected curve surjects onto that of the curve, as in the proof of Theorem \bpout{14.26}, and the same holds for the Weil groups. Hence \(Q\) or \(R\) has zero generic restriction. Such a perverse object is punctual. Indeed, its \(\mathcal H^{-1}\) is then supported on a finite set, and a sheaf without punctual subsheaves with this property is zero, while its \(\mathcal H^0\) is punctual by definition. Thus \(Q\) or \(R\) is a punctual subobject or quotient of \(P\), which by the vanishing just proved is zero, and so \(P\) is simple.

Conversely, let \(P\) be simple. We show that \(P\) is punctual or a middle extension of an irreducible lisse sheaf. Ordinary truncation gives a perverse exact sequence

\[
0\longrightarrow(\mathcal H^{-1}P)[1]
\longrightarrow P\longrightarrow\mathcal H^0P\longrightarrow0.
\]

Hence a simple \(P\) is either punctual or of the form \(F[1]\), where \(F\) has no punctual subsheaf. In the second case, choose a dense open \(U\) on which \(F\) is lisse and put \(M=F|_U\). Then the map \(F\to j_*M\) is injective, since its kernel is \(F_{\rm p}=0\) (Lemma \hyperref[bp33]{16.10}), and its cokernel \(B\) is punctual. Rotating the corresponding triangle of sheaves we obtain the perverse exact sequence \(0\to B\to F[1]\to(j_*M)[1]\to0\), whose last term is nonzero, so \(B=0\) by simplicity. It remains to see that \(M\) is irreducible. Suppose that \(M\) had a proper nonzero lisse subsheaf \(M_1\). Then the object \(j_!M_1[1]\) is perverse, because \(j_!M_1\subset j_*M_1\) has no punctual subsheaf, and the image of \(j_!M_1[1]\to P\) in the abelian heart is proper and nonzero, as one sees on \(U\). This contradicts simplicity. Finally, if \(P\) is punctual, then by the long exact cohomology sequence its subobjects in the heart are punctual sheaves in degree zero, so \(P\) is simple exactly when it is irreducible as a punctual sheaf.
\end{proof}

% object: example-extension-by-zero-not-simple (blueprint/chap-The-Fourier-transform/sec-Perverse-sheaves-on-a-curve/example-extension-by-zero-not-simple/_section.md)
\phantomsection\label{bp45}%
\begin{bpstatement}{Example~16.15}{ (Extension by zero is not simple)}
Let \(j:\mathbf G_m\hookrightarrow\mathbf A^1\) and \(i:0\hookrightarrow\mathbf A^1\) be the inclusions over \(k\), and put \(\delta_0=i_*E\). The sheaf \(j_!E\) has no punctual subsheaf, as its stalk at the origin is zero, so \(j_!E[1]\) is perverse. It is not simple. Rotating the triangle of the exact sequence \(0\to j_!E\to E\to\delta_0\to0\) gives the triangle \(\delta_0\to j_!E[1]\to E[1]\to\). Its three terms are perverse, and so

\[
0\longrightarrow\delta_0\longrightarrow j_!E[1]
\longrightarrow E[1]\longrightarrow0
\]

is exact in the perverse heart. Thus, while among sheaves the skyscraper is a quotient of \(E\), in the perverse heart it becomes a subobject of \(j_!E[1]\). By contrast, the middle extension of \(E|_{\mathbf G_m}\) is \(E\) itself (Lemma \bpout{15.2}), and \(E[1]\) is simple. On the other hand \((E\oplus\delta_0)[1]\) is not perverse, because \(\mathcal H^{-1}\) has the punctual subsheaf \(\delta_0\), which is the subsheaf \(F_{\rm p}\) of Lemma \hyperref[bp33]{16.10}, the kernel of \(E\oplus\delta_0\to j_*j^*(E\oplus\delta_0)=E\).
\end{bpstatement}


% object: sec-Perverse-t-exactness (blueprint/chap-The-Fourier-transform/sec-Perverse-t-exactness/_section.md), content hidden
\subsection*{\texorpdfstring{16.5\quad Perverse t-exactness}{16.5 Perverse t-exactness}}
\phantomsection\label{bp46}%
\addcontentsline{toc}{subsection}{\texorpdfstring{16.5\quad Perverse t-exactness}{16.5 Perverse t-exactness}}
% text: prose-16-15
\phantomsection\label{bp47}%
The transform \(\operatorname{FT}_\psi\) is t-exact for the perverse t-structure on the line (Proposition \hyperref[bp48]{16.16}). By the stalk formula (16.2), this comes down to the compactly supported cohomology of sheaves on an affine curve in the extreme degrees zero and two.

% object: proposition-fourier-transform-perverse-heart (blueprint/chap-The-Fourier-transform/sec-Perverse-t-exactness/proposition-fourier-transform-perverse-heart/_section.md)
\phantomsection\label{bp48}%
\begin{bpstatementit}{Proposition~16.16}{ (Fourier transform and the perverse heart)}
On the geometric affine line, \(\operatorname{FT}_\psi\) is a t-exact equivalence of bounded constructible derived categories for the perverse t-structure of Lemma \hyperref[bp37]{16.12}, with inverse \(\operatorname{FT}_{\psi^{-1}}(-)(1)\). Hence it is an exact equivalence of the perverse hearts, and its canonical Frobenius compatibilities induce an exact equivalence of the categories of perverse Weil sheaves.
\end{bpstatementit}

% proof: proof/proposition-fourier-transform-perverse-heart
\begin{proof}[{Proof}]\phantomsection\label{bp49}%
We show that the geometric transform and its inverse preserve \({}^pD^{\le0}\). Being mutually inverse, they then preserve \({}^pD^{\ge0}\) by orthogonality, and the Weil statement follows by Lemma \hyperref[bp40]{16.13}, since the forgetful functor on perverse Weil sheaves is exact and reflects exactness. For the first step, the key point is that the stalks of the transform are compactly supported cohomology groups on the affine line, in degrees \(0,1,2\), and that the group in degree two is nonzero only at the finitely many points \(y\) for which the twisted sheaf has a geometrically constant quotient.

\emph{Right exactness.} Since the transform is triangulated and \({}^pD^{\le0}\) is stable under extensions and positive shifts, ordinary truncation reduces this to punctual sheaves in degree zero and to objects \(G[1]\) with \(G\) constructible. Indeed, an object of \({}^pD^{\le0}\) is an iterated extension of its shifted ordinary cohomology sheaves \(\mathcal H^a(K)[-a]\) with \(a\le0\); the one with \(a=0\) is punctual, and the others are of the form \(G[1]\) shifted by a nonnegative amount. Punctual sheaves map into \({}^pD^{\le0}\), since by the stalk formula the transform of a geometric skyscraper is a lisse sheaf shifted by one.

For \(G[1]\) we show that the transform lies in ordinary degrees \(-2,-1,0\) and that its degree-zero sheaf has finite support. The first assertion follows from the stalk formula and Lemma \bpout{15.8}. For the second, note that the degree-zero stalk at a closed geometric point \(y\) is

\[
H_c^2(\mathbf A^1,G\otimes\mathcal L_{\psi(yx)}).
\]

On a dense open where \(G\) is lisse with generic representation \(V\), write \(\chi_y\) for the character of \(\pi_1\) given by the restriction of \(\mathcal L_{\psi(yx)}\). By localization we may discard the finitely many remaining points, which contribute only in degree zero, and Lemma \bpout{15.7} (whose proof of the formula for \(H^2_c\) uses only the geometric curve) shows that this stalk is \((V\otimes\chi_y)_{\pi_1}(-1)\). In particular, it is nonzero only if \(V\) has the rank-one quotient \(\chi_y^{-1}\), and we show that this happens for only finitely many \(y\). First, the characters \(\chi_y\) are pairwise distinct. Indeed, if \(\chi_y=\chi_{y'}\) with \(y\ne y'\), then \(\mathcal L_{\psi((y-y')x)}\) is trivial on the dense open, hence on the line by the surjectivity of fundamental groups. This is impossible, because its compactly supported cohomology vanishes by Lemma \hyperref[bp13]{16.4}, whereas the constant sheaf has nonzero top cohomology with compact support. Second, if \(V\) has pairwise nonisomorphic rank-one quotients \(\chi_1,\dots,\chi_n\), then the induced map \(V\to\chi_1\oplus\dots\oplus\chi_n\) is surjective. To see this, note that its image is a subrepresentation of a sum of pairwise nonisomorphic simple representations, hence the sum of some of the \(\chi_i\), and it maps onto each \(\chi_i\). So \(n\le\dim V\), and \(\chi_y^{-1}\) is a quotient of \(V\) for at most \(\dim V\) points \(y\). The degree-zero sheaf therefore has only finitely many nonzero closed geometric stalks. Since it is constructible, it has finite support, and the transform of \(G[1]\) belongs to \({}^pD^{\le0}\).

\emph{Exactness.} The argument for \(\operatorname{FT}_\psi\) applies equally to \(\psi^{-1}\), and a Tate twist does not affect supports or degrees. Both \(\operatorname{FT}_\psi\) and its quasi-inverse \(\operatorname{FT}_{\psi^{-1}}(-)(1)\) of Proposition \hyperref[bp19]{16.6} therefore preserve \({}^pD^{\le0}\). Since \({}^pD^{\ge0}\) is the right orthogonal of \({}^pD^{\le-1}\), the two mutually inverse equivalences also preserve \({}^pD^{\ge0}\). Both functors are therefore t-exact, and they induce exact equivalences of the geometric hearts. With their canonical Frobenius compatibilities and the equivariant inversion isomorphisms of Proposition \hyperref[bp19]{16.6}, they then give quasi-inverse functors on perverse Weil sheaves, which are exact because the forgetful functors are exact and reflect exactness.
\end{proof}


% object: sec-Middle-extensions (blueprint/chap-The-Fourier-transform/sec-Middle-extensions/_section.md), content hidden
\subsection*{\texorpdfstring{16.6\quad Middle extensions}{16.6 Middle extensions}}
\phantomsection\label{bp50}%
\addcontentsline{toc}{subsection}{\texorpdfstring{16.6\quad Middle extensions}{16.6 Middle extensions}}
% text: prose-16-16
\phantomsection\label{bp51}%
Let \(j:U_0\hookrightarrow\mathbf A^1_{k_0}\) be a dense open, and let \(M\) be an irreducible lisse Weil sheaf on \(U_0\). The argument of Chapter \bpout{17} needs the transform of the middle extension of \(M\) to be again a middle extension, that is, to come from a lisse sheaf \(M'\) on a dense open of the dual line. This holds unless \(M\) has a geometric constituent of the form \(\mathcal L_{\psi(ax)}\), and the hypothesis of the next proposition excludes these. They behave differently: for \(a\ne0\) the sheaf \(\mathcal L_{\psi(ax)}\) has no cohomology with compact support on the line by Lemma \hyperref[bp13]{16.4}, and its transform is a skyscraper rather than a lisse sheaf of weight one. Excluding them, however, costs nothing in the application, because a sheaf unramified at infinity has no such constituent other than the constant one, and constant constituents contribute nothing, since \(H^1_c(\mathbf A^1,\overline{\mathbf Q}_\ell)=0\).

% object: proposition-fourier-transform-middle-extension (blueprint/chap-The-Fourier-transform/sec-Middle-extensions/proposition-fourier-transform-middle-extension/_section.md)
\phantomsection\label{bp52}%
\begin{bpstatementit}{Proposition~16.17}{ (The Fourier transform of a middle extension)}
Suppose that \(M\) is nonzero and irreducible, and that none of its geometric irreducible constituents is the restriction of \(\mathcal L_{\psi(ax)}\) for some \(a\in\bar k_0\) (the case \(a=0\) being the constant sheaf). Then there are a dense open \(j':U'_0\hookrightarrow\mathbf A^{1\vee}_{k_0}\) and a nonzero irreducible lisse Weil sheaf \(M'\) on it such that

\[
\operatorname{FT}_\psi((j_*M)[1])\simeq(j'_*M')[1].
\]
\end{bpstatementit}

% proof: proof/proposition-fourier-transform-middle-extension
\begin{proof}[{Proof}]\phantomsection\label{bp53}%
By Fourier inversion, a punctual transform would force every geometric constituent of \(M\) to be a linear character sheaf, and the hypothesis excludes this.

By Proposition \hyperref[bp48]{16.16}, the transform of the simple object \((j_*M)[1]\) is simple, and so by Lemma \hyperref[bp43]{16.14} it is either punctual or of the displayed form. Restricting the kernel to \(\{a\}\times\mathbf A^1\) shows that the inverse transform of a geometric skyscraper at \(a\) is \(\mathcal L_{\psi(-ax)}(1)[1]\). Now, if the transform were punctual, it would geometrically be an iterated extension of skyscrapers at finitely many points \(a\), so by Fourier inversion and exactness \((j_*M)[1]\) would be geometrically an iterated extension of the objects \(\mathcal L_{\psi(-ax)}(1)[1]\). Restricting to \(U\), we see that every geometric constituent of \(M\) would be one of these linear characters, which the hypothesis excludes. Finally, the open \(U'_0\) and the sheaf \(M'\) are defined over \(k_0\), since we may choose a Weil-stable open on which the constructible transform is lisse.
\end{proof}

% text: prose-16-17
\phantomsection\label{bp54}%
The sheaf \(M'\) is determined by \(M\) up to isomorphism and shrinking of \(U'_0\), since it is the restriction to \(U'_0\) of \(\mathcal H^{-1}\operatorname{FT}_\psi((j_*M)[1])\). In the application we check the hypothesis of the proposition with the following criterion, whose proof is a computation of ramification at infinity.

% object: lemma-unramified-at-infinity-no-linear-characters (blueprint/chap-The-Fourier-transform/sec-Middle-extensions/lemma-unramified-at-infinity-no-linear-characters/_section.md)
\phantomsection\label{bp55}%
\begin{bpstatementit}{Lemma~16.18}{ (Sheaves unramified at infinity)}
Let \(M\) be an irreducible lisse Weil sheaf on a dense open \(U_0\subset\mathbf A^1_{k_0}\). If \(M\) is unramified at infinity and is not geometrically constant, then no geometric constituent of \(M\) is of the form \(\mathcal L_{\psi(ax)}\) with \(a\in\bar k_0\).
\end{bpstatementit}

% proof: proof/lemma-unramified-at-infinity-no-linear-characters
\begin{proof}[{Proof}]\phantomsection\label{bp56}%
For \(a\ne0\) the sheaf \(\mathcal L_{\psi(ax)}\) is ramified at infinity, and the trivial character is excluded because \(M\) is irreducible and not geometrically constant, its geometric restriction being semisimple. For the first point, write \(t=1/x\) and express the character by the Artin–Schreier equation \(u^p-u=c/t\) with \(c\ne0\). Let \(v\) be the normalized valuation of the field generated by a root \(u\) over the completion \(k((t))\) at infinity, and let \(e\) be the ramification index, so that \(v(t)=e\) and \(e\le p\). Since \(v(c/t)=-e<0\), we must have \(v(u)<0\), so that \(v(u^p)=pv(u)<v(u)\) and \(v(u^p-u)=pv(u)\). Now \(pv(u)=-e\) forces \(p\mid e\), and so \(e=p\): the cover is totally ramified. Inertia at infinity therefore maps onto the Galois group \(\mathbf F_p\) of this cover, on which \(\psi\) is nontrivial. Since \(M\) is unramified at infinity, this proves the first point, because a character that is nontrivial on inertia cannot be a constituent of a representation on which inertia acts trivially.

For the trivial character, we apply Lemma \bpout{14.11} to the geometric fundamental group, which is normal in the Weil group: the restriction of \(M\) to it is semisimple, and so a trivial geometric constituent gives nonzero geometric invariants. Since this space of invariants is Weil-stable, the irreducibility of \(M\) shows that \(M\) is geometrically constant.
\end{proof}



\clearpage\phantomsection\hypertarget{bp-outside}{}%
\section*{Outside this export}
\addcontentsline{toc}{section}{Outside this export}
References marked \ensuremath{{}^\ast} are to parts of the blueprint this document leaves out:
\begin{description}
\item[Definition 9.24] Trace map, in Ⅱ Duality and the trace formula (\bpslug{definition-{}\allowbreak{}trace-{}\allowbreak{}map})
\item[Definition 10.18] Weil structure, in Ⅱ Duality and the trace formula (\bpslug{definition-{}\allowbreak{}weil-{}\allowbreak{}structure})
\item[Proposition 11.2] Properties of the trace, in Ⅱ Duality and the trace formula (\bpslug{proposition-{}\allowbreak{}properties-{}\allowbreak{}trace})
\item[Section 12.6] The function–sheaf dictionary, in Ⅱ Duality and the trace formula (\bpslug{sec-{}\allowbreak{}The-{}\allowbreak{}function-{}\allowbreak{}sheaf-{}\allowbreak{}dictionary-{}\allowbreak{}2})
\item[Proposition 12.28] The function–sheaf dictionary, in Ⅱ Duality and the trace formula (\bpslug{proposition-{}\allowbreak{}function-{}\allowbreak{}sheaf-{}\allowbreak{}dictionary})
\item[Proposition 13.10] Lattices and representations on a normal scheme, in Ⅲ Weights and the Riemann hypothesis (\bpslug{proposition-{}\allowbreak{}lattices-{}\allowbreak{}representations-{}\allowbreak{}normal-{}\allowbreak{}scheme})
\item[Lemma 14.11] Restriction to a normal subgroup, in Ⅲ Weights and the Riemann hypothesis (\bpslug{lemma-{}\allowbreak{}restriction-{}\allowbreak{}normal-{}\allowbreak{}subgroup})
\item[Theorem 14.26] Purity of real sheaves on curves, in Ⅲ Weights and the Riemann hypothesis (\bpslug{theorem-{}\allowbreak{}purity-{}\allowbreak{}real-{}\allowbreak{}sheaves-{}\allowbreak{}curves})
\item[Lemma 15.2] The boundary stalk, in Ⅲ Weights and the Riemann hypothesis (\bpslug{lemma-{}\allowbreak{}boundary-{}\allowbreak{}stalk})
\item[Lemma 15.7] The extreme degrees on an affine curve, in Ⅲ Weights and the Riemann hypothesis (\bpslug{lemma-{}\allowbreak{}extreme-{}\allowbreak{}degrees-{}\allowbreak{}affine-{}\allowbreak{}curve})
\item[Lemma 15.8] Cohomological range on a curve, in Ⅲ Weights and the Riemann hypothesis (\bpslug{lemma-{}\allowbreak{}cohomological-{}\allowbreak{}range-{}\allowbreak{}curve})
\item[Definition 15.10] Growth of squared traces, in Ⅲ Weights and the Riemann hypothesis (\bpslug{definition-{}\allowbreak{}growth-{}\allowbreak{}squared-{}\allowbreak{}traces})
\item[Theorem 15.14] Weights and sums of squared traces, in Ⅲ Weights and the Riemann hypothesis (\bpslug{theorem-{}\allowbreak{}weights-{}\allowbreak{}sums-{}\allowbreak{}squared-{}\allowbreak{}traces})
\item[Chapter 17] The Deligne–Laumon theorem for curves, in Ⅲ Weights and the Riemann hypothesis (\bpslug{chap-{}\allowbreak{}The-{}\allowbreak{}Deligne-{}\allowbreak{}Laumon-{}\allowbreak{}theorem-{}\allowbreak{}for-{}\allowbreak{}curves})
\item[Appendix D] Cohomology of torsion sheaves on curves, in Appendices (\bpslug{chap-{}\allowbreak{}Cohomology-{}\allowbreak{}of-{}\allowbreak{}torsion-{}\allowbreak{}sheaves-{}\allowbreak{}on-{}\allowbreak{}curves})
\end{description}
\clearpage\phantomsection%
\section*{References}
\addcontentsline{toc}{section}{References}
\bpbibitem{BBD82}{bp58}{A. Beilinson, J. Bernstein, and P. Deligne. \emph{Faisceaux pervers}, volume 100 of \emph{Astérisque}. 1982.}

\bpbibitem{BS15}{bp59}{B. Bhatt and P. Scholze. The pro-étale topology for schemes. \emph{Astérisque}, 369:99–201, 2015. \href{https://people.mpim-bonn.mpg.de/scholze/proetale.pdf}{Author's manuscript}.}

\bpbibitem{Con07}{bp60}{B. Conrad. Deligne's notes on Nagata compactifications. \emph{Journal of the Ramanujan Mathematical Society}, 22:205–257, 2007. \href{https://math.stanford.edu/\~conrad/papers/nagatafinal.pdf}{Author's manuscript}.}

\bpbibitem{Del74}{bp61}{P. Deligne. La conjecture de Weil. I. \emph{Publ. Math. IHÉS}, 43:273–307, 1974.}

\bpbibitem{Del80}{bp62}{P. Deligne. La conjecture de Weil. II. \emph{Publ. Math. IHÉS}, 52:137–252, 1980.}

\bpbibitem{Den22}{bp63}{C. Deninger. There is no “Weil-”cohomology theory with real coefficients for arithmetic curves. \emph{Ann. Sc. Norm. Super. Pisa Cl. Sci.}, 2022. To appear. arXiv:\href{https://arxiv.org/abs/2204.02714}{2204.02714}.}

\bpbibitem{Deu41}{bp64}{M. Deuring. Die Typen der Multiplikatorenringe elliptischer Funktionenkörper. \emph{Abh. Math. Sem. Univ. Hamburg}, 14:197–272, 1941.}

\bpbibitem{Eke90}{bp65}{T. Ekedahl. On the adic formalism. In \emph{The Grothendieck Festschrift, Vol. II}, volume 87 of \emph{Progress in Mathematics}, pages 197–218. Birkhäuser, 1990.}

\bpbibitem{Has36}{bp66}{H. Hasse. Zur Theorie der abstrakten elliptischen Funktionenkörper. I, II, III. \emph{J. Reine Angew. Math.}, 175:55–62, 69–88, 193–208, 1936.}

\bpbibitem{Hum72}{bp67}{J. E. Humphreys. \emph{Introduction to Lie Algebras and Representation Theory}, volume 9 of \emph{Graduate Texts in Mathematics}. Springer, 1972.}

\bpbibitem{Ill94}{bp68}{L. Illusie. Autour du théorème de monodromie locale. \emph{Astérisque}, 223:9–57, 1994. §4. \href{https://www.numdam.org/item/AST_1994__223__9_0/}{Original article}.}

\bpbibitem{Kat88}{bp69}{N. M. Katz. Travaux de Laumon. \emph{Astérisque}, 161–162:105–132, 1988. Séminaire Bourbaki, exp. 691.}

\bpbibitem{KW01}{bp70}{R. Kiehl and R. Weissauer. \emph{Weil conjectures, perverse sheaves and \(\ell\)-adic Fourier transform}, volume 42 of \emph{Ergebnisse der Mathematik}. Springer, 2001.}

\bpbibitem{Lau87}{bp57}{G. Laumon. Transformation de Fourier, constantes d'équations fonctionnelles et conjecture de Weil. \emph{Publ. Math. IHÉS}, 65:131–210, 1987.}

\bpbibitem{Mil86a}{bp71}{J. S. Milne. \href{https://www.jmilne.org/math/xnotes/AVs.pdf}{Abelian Varieties}. In \emph{Arithmetic Geometry}. 1986. Corrected online version.}

\bpbibitem{Mil86b}{bp72}{J. S. Milne. \href{https://www.jmilne.org/math/xnotes/JVs.pdf}{Jacobian Varieties}. In \emph{Arithmetic Geometry}. 1986. Corrected online version.}

\bpbibitem{Ran39}{bp73}{R. A. Rankin. Contributions to the theory of Ramanujan's function \(\tau(n)\) and similar arithmetical functions. II. \emph{Proc. Cambridge Philos. Soc.}, 35:357–372, 1939.}

\bpbibitem{SGA4}{bp74}{M. Artin, A. Grothendieck, and J.-L. Verdier, editors. \emph{Théorie des topos et cohomologie étale des schémas}, volumes 269, 270, 305 of \emph{Lecture Notes in Mathematics}. Springer, 1972–1973. SGA 4.}

\bpbibitem{SGA4½}{bp75}{P. Deligne. \emph{Cohomologie étale}, volume 569 of \emph{Lecture Notes in Mathematics}. Springer, 1977. SGA 4½; \emph{Théorèmes de finitude}, §3, and \emph{Rapport sur la formule des traces}, §3, and \emph{Sommes trigonométriques}, §1.}

\bpbibitem{SGA5}{bp76}{A. Grothendieck, editor. \emph{Cohomologie \(\ell\)-adique et fonctions \(L\)}, volume 589 of \emph{Lecture Notes in Mathematics}. Springer, 1977. SGA 5.}

\bpbibitem{SGA7II}{bp77}{P. Deligne and N. Katz, editors. \emph{Groupes de monodromie en géométrie algébrique}, volume 340 of \emph{Lecture Notes in Mathematics}. Springer, 1973. SGA 7 II; Exposé XIII, (2.1.7).}

\bpbibitem{Stacks}{bp78}{The Stacks Project Authors. \href{https://stacks.math.columbia.edu}{The Stacks Project}. Individual results are cited by their stable tags in the text.}

\bpbibitem{Var05}{bp79}{Y. Varshavsky. Lefschetz–Verdier trace formula and a generalization of a theorem of Fujiwara. 2005. Version 2. arXiv:\href{https://arxiv.org/abs/math/0505564}{math/0505564}. \href{https://arxiv.org/abs/math/0505564v2}{Author's preprint}.}

\bpbibitem{Wei48a}{bp80}{A. Weil. \emph{Sur les courbes algébriques et les variétés qui s'en déduisent}. Hermann, Paris, 1948.}

\bpbibitem{Wei48b}{bp81}{A. Weil. \emph{Variétés abéliennes et courbes algébriques}. Hermann, Paris, 1948.}

\bpbibitem{Wei49}{bp82}{A. Weil. Numbers of solutions of equations in finite fields. \emph{Bull. Amer. Math. Soc.}, 55:497–508, 1949.}

\end{document}
